\documentclass[12pt,openany]{memoir}

\def\nsemestre{II}
\def\nterm{Fall}
\def\nyear{2026}
\def\nprofesor{Barbara Fantechi}
\def\nauthor{Ignacio Rojas}
\def\nsigla{MMARGS}
\def\nsiglahead{Stacks: de las familias a los espacios de módulos}
% No definimos \nlang: así el encabezado compartido activa el español.
\PassOptionsToPackage{es-noshorthands,es-noquoting}{babel}
\let\footruleskip\relax
% This document is one directory deeper than Stacks4StupidPeople.tex.

\makeatletter
\ifx \nauthor\undefined
  \def\nauthor{Ignacio Rojas}
\else
\fi

\ifx \nextra \undefined
\ifx \nlang \undefined
\author{Basado en las clases impartidas por \nprofesor \\\small Notas tomadas por \nauthor}
\else
\author{Based on the lectures by \nprofesor \\\small Notes written by \nauthor}
\fi
\else
\author{\nauthor}
\fi
\date{\nterm\ \nyear}

%%%%%%%%%%%%%
%% 1. Pacotes
%%%%%%%%%%%%%

\usepackage{alltt}
\usepackage{amsfonts}
\usepackage{amsmath}
\usepackage{amssymb}
\usepackage{amsthm}
\usepackage{algorithm}
\usepackage[noend]{algpseudocode}
\usepackage{array}
\newcommand\hmmax{0} % default 3
\newcommand\bmmax{0} % default 4 %%tex.se/3676,219310
%\usepackage{bbold}
\usepackage{bm}
\usepackage{booktabs}
%\usepackage{caption}
%\usepackage{cancel}
%\usepackage{dsfont}
\usepackage{esint}
\usepackage{fancyhdr}
\usepackage{graphicx}
\usepackage[utf8]{inputenc}
\usepackage{listings}
\usepackage{mathabx}
\usepackage[cal=euler]{mathalfa}
%\usepackage[cal=euler,frak=euler]{mathalfa} % mathcal (JIRR) precisabamos correr initexmf --mkmaps en cmd JCVDG
\usepackage{mathdots}
\usepackage{mathrsfs}
%\usepackage{mathtools}
\usepackage{microtype}
\usepackage{multicol}
\usepackage{multirow}
\usepackage[theoremfont,largesc,tighter,osf]{newpxtext} %JCV Diff
\let\widering\undefined
%\usepackage[bigdelims,vvarbb]{newpxmath} %JCVDG
%por alguna razón esto afectaba las tildes en \min, \lim y demás
%\usepackage{pdflscape}
\usepackage{pgfplots}
\usepackage{physics}
\usepackage{siunitx}
\usepackage{slashed}
\usepackage{stmaryrd}
\SetSymbolFont{stmry}{bold}{U}{stmry}{m}{n}
%\usepackage{subfigure}
\usepackage{subcaption}
\usepackage{tabularx}
\usepackage[breakable,skins]{tcolorbox}
\usepackage{textcomp} %%JCVDG
\usepackage{tikz}
\usepackage{tkz-euclide}
\usepackage[normalem]{ulem}
\usepackage[all]{xy}
\usepackage{imakeidx}
\ifx \nlang \undefined
\usepackage[spanish]{babel}
\else\fi 
\usepackage{wrapfig}

%%%%%%%%%%%%%%%%%%%%
%% 2. Document Setup
%%%%%%%%%%%%%%%%%%%%

\ifx \nextra \undefined
    \ifx \nlang \undefined
    \makeindex[intoc, title=Índice Analítico] %Título de índice analítico
    %El índice general es aquel en el que se indican los capítulos, títulos y subtítulos del libro.
    %Índice onomástico es donde aparece el nombre de personas mencionadas en el texto, por orden alfabético con el número de las páginas donde aparecen.
    %El índice analítico se refiere a los temas y conceptos que aparecen en el libro
    \indexsetup{othercode={\fancyhead[LE]{\emph{Índice Analítico}}}}
    \else
    \makeindex[intoc, title=Index] 
    \indexsetup{othercode={\fancyhead[LE]{\emph{Index}}}}
    \fi
  \usepackage[pdftex,
    hidelinks,
    pdfauthor={\nauthor},
    pdfsubject={Notas: \nsiglahead\ \nsemestre-\nyear},
    pdftitle={Semestre \nsemestre\ - \nsigla},
  pdfkeywords={UCR Costa Rica Matem\'aticas Mate \nsemestre\ \nterm\ \nyear\ \nsiglahead}]{hyperref}
  \title{\nsigla\ --- \nsiglahead}
\else
  \usepackage[pdftex,
     hidelinks,
    pdfauthor={\nauthor},
    pdfsubject={\nextra \nsiglahead\ \nsemestre-\nyear},
    pdftitle={Semestre \nsemestre\ - \nsigla},
  pdfkeywords={UCR Costa Rica Matem\'aticas Mate \nsemestre\ \nterm\ \nyear\ \nsiglahead\ \nextra}]{hyperref}

  \title{\nsigla\ --- \nsiglahead \\ {\Large \nextra}}
  \renewcommand\printindex{}
\fi

\pgfplotsset{compat=1.12}


\pagestyle{fancy}
\setlength{\headheight}{15.72pt} %preceding warning said make it at least this


\ifx \nsiglahead \undefined
\def\nsiglahead{\nsigla}
\fi

\lhead{} %%%empty lhead
\rfoot{\thepage}

\ifx \nextra \undefined
  \chead{
    \ifnum\thepage=1
    \else
      \ifx \nlang \undefined
      \textbf{Notas \nsiglahead\ \nsemestre-\nyear}
      \else
      \textbf{Notes \nsiglahead\ \nsemestre-\nyear}
      \fi
    \fi}
  \rhead{}%\firstxmark} % Top right header
\else
%    \chead{
%    \ifnum\thepage=1
%    \else
%      \textbf{Notas \nsiglahead\ \nsemestre-\nyear \ (\nextra)}
%    \fi}
     \chead{
       \textbf{\nextra\ \nsigla\ \nsemestre-\nyear}
     }
     \rhead{
       \textbf{\nauthor}
     }
\fi
\lfoot{}%\lastxmark} % Bottom left footer
\cfoot{} % Bottom center footer

\usetikzlibrary{arrows.meta}
\usetikzlibrary{decorations.markings}
\usetikzlibrary{decorations.pathmorphing}
\usetikzlibrary{positioning}
\usetikzlibrary{fadings}
\usetikzlibrary{intersections}
\usetikzlibrary{cd}

\ifx \nhtml \undefined
\else
  \renewcommand\printindex{}
  \DisableLigatures[f]{family = *}
  \let\Contentsline\contentsline
  \renewcommand\contentsline[3]{\Contentsline{#1}{#2}{}}
  \renewcommand{\@dotsep}{10000}
  \newlength\currentparindent
  \setlength\currentparindent\parindent

  \newcommand\@minipagerestore{\setlength{\parindent}{\currentparindent}}
  \usepackage[active,tightpage,pdftex]{preview}
  \renewcommand{\PreviewBorder}{0.1cm}

  \newenvironment{stretchpage}%
  {\begin{preview}\begin{minipage}{\hsize}}%
    {\end{minipage}\end{preview}}
  \AtBeginDocument{\begin{stretchpage}}
  \AtEndDocument{\end{stretchpage}}

  \newcommand{\@@newpage}{\end{stretchpage}\begin{stretchpage}}

  \let\@real@section\section
  \renewcommand{\section}{\@@newpage\@real@section}
  \let\@real@subsection\subsection
  \renewcommand{\subsection}{\@ifstar{\@real@subsection*}{\@@newpage\@real@subsection}}
\fi
\ifx \ntrim \undefined
\usepackage[shortlabels]{enumitem} %mfw package order matters por savetrees
\else
  \usepackage{geometry}
  \geometry{
    papersize={379pt, 699pt},
    textwidth=345pt,
    textheight=596pt,
    left=17pt,
    top=54pt,
    right=17pt
  }
  \headwidth=345pt
 \usepackage[extreme]{savetrees}
\fi

\ifx \darktheme\undefined
\else
\pagecolor[rgb]{0,0,0}
\color[rgb]{0.5,0.5,0.5}
%\pagecolor[rgb]{0.2,0.231,0.302}%{0.23,0.258,0.321}
%\color[rgb]{1,1,1}
\fi

\ifx\nthesis \undefined
\ifx \nextra \undefined
\let\@real@maketitle\maketitle
\renewcommand{\maketitle}{\@real@maketitle\begin{center}\begin{minipage}[c]{0.9\textwidth}\centering\footnotesize 
  \ifx \nlang \undefined
  Estas notas no están respaldadas por los profesores y han sido modificadas (a menudo de manera significativa) después de las clases. No están lejos de ser representaciones precisas de lo que realmente se dio en clase y en particular todos los errores son casi seguramente míos.
  \else 
  Please note that these notes were not provided or endorsed by the lecturer and have been significantly altered after the class. They may not accurately reflect the content covered in class and any errors are solely my responsibility.
  \fi
\end{minipage}\end{center}}
%\else decia que habia un extra
\else
\fi
\fi

\def\moverlay{\mathpalette\mov@rlay}
\def\mov@rlay#1#2{\leavevmode\vtop{%
   \baselineskip\z@skip \lineskiplimit-\maxdimen
   \ialign{\hfil$\m@th#1##$\hfil\cr#2\crcr}}}
\newcommand{\charfusion}[3][\mathord]{
    #1{\ifx#1\mathop\vphantom{#2}\fi
        \mathpalette\mov@rlay{#2\cr#3}
      }
    \ifx#1\mathop\expandafter\displaylimits\fi}

%%%%%%%%%%%%%%%%%%%%%%%%%%%%%%
%% 2.1 Some internal machinery
%%%%%%%%%%%%%%%%%%%%%%%%%%%%%%

\makeatletter
\renewcommand{\section}{\@startsection{section}{1}{\z@}%
							 {-3.25ex \@plus -1ex \@minus -.2ex}%
							 {1.5ex \@plus.2ex}%
							 {\normalfont\large\bfseries}}
\renewcommand{\subsection}{\@startsection{subsection}{2}{\z@}%
							 {-3.25ex \@plus -1ex \@minus -.2ex}%
							 {1.5ex \@plus .2ex}%
               {\normalfont\normalsize\bfseries}}
\newcommand*{\defeq}{\!\mathrel{\rlap{%
             \raisebox{0.3ex}{$\m@th\cdot$}}%
             \raisebox{-0.3ex}{$\m@th\cdot$}}%
                    =\!}
\makeatother
\ifx\ntrim\undefined
\newcommand{\coursetitle}{\nsigla: \nsiglahead}
\ifx\nextra\undefined
\pagestyle{ruled}
\makeoddhead{ruled}{\coursetitle}{}{\rightmark}
\else\fi
\settypeblocksize{49pc}{37pc}{*}
\setlrmargins{*}{*}{1.2}
\setulmargins{*}{*}{0.8}
\setheadfoot{16pt}{30pt}
\setheaderspaces{*}{1.5pc}{1}
\setmarginnotes{1pt}{1pt}{1pt}
\checkandfixthelayout

\setlength{\unitlength}{3pt}
\setlength{\hfuzz}{1pt}

\setlength{\fboxsep}{6pt}

\setlength{\footskip}{17pt}

\linespread{1.1}
\else\fi
\renewcommand{\cftdotsep}{\cftnodots} %%% no dots in ToC
\setpnumwidth{2em}  %%% width of page-number box in ToC


\newcommand{\stophere}{\relax} %% can be changed to `\endinput'
% \newcommand{\stophere}{\endinput} %% can be changed to `\relax'


\DeclareRobustCommand{\qned}{\ifmmode
  \else \leavevmode\unskip\penalty9999 \hbox{}\nobreak\hfill \fi
  \quad\hbox{\qnedsymbol}}
\newcommand{\qnedsymbol}{$\boxminus$} %% No-proofs end with `\qned'

\DeclareRobustCommand{\qef}{\ifmmode
  \else \leavevmode\unskip\penalty9999 \hbox{}\nobreak\hfill \fi
  \quad\hbox{\qefsymbol}}
\newcommand{\qefsymbol}{$\lozenge$} %% Examples end with `\qef'
\def\enddefn{\qef\endtrivlist}      %% `\qef' automático en defns
\def\endejem{\qef\endtrivlist}      %% `\qef' automático en ejemplos

\newcommand{\hideqed}{\renewcommand{\qed}{}} %% to suppress `\qed'
\newcommand{\hideqef}{\renewcommand{\qef}{}} %% to suppress `\qef'

% \newcommand{\ldbrack}{\ensuremath{[\mskip-2.5mu[}} %% corchetes [[
% \newcommand{\rdbrack}{\ensuremath{]\mskip-2.5mu]}} %% corchetes ]]

\newcommand{\stroke}{\mathbin|}     %% (for `\bbraket' and such)

\newcommand{\rtri}{\blacktriangleright} %% (for `\marker' and such)
\newcommand{\tribar}{|\mkern-2mu|\mkern-2mu|} %% norma triple: |||


%% Formatting changes:

\renewcommand{\labelitemi}{$\diamond$} %% instead of bullets

\renewcommand{\theenumi}{\alph{enumi}}  %% use lowercase letters
\renewcommand{\labelenumi}{\textup{(\theenumi)}} %% inside parentheses

%%%%%%%%%%%%%%
%% 2.2. Colors
%%%%%%%%%%%%%%

\definecolor{MATLABgreen}{RGB}{28,172,0} % color values Red, Green, Blue
\definecolor{MATLABlila}{RGB}{170,55,241}
\definecolor{dankBlue}{RGB}{51,60,77} % color values Red, Green, Blue
\definecolor{dankBlueLite}{RGB}{82,97,125} % color values Red, Green, Blue
\definecolor{celesUCR}{RGB}{0,192,243}
\definecolor{azulUCR}{RGB}{0,93,164}
\definecolor{verdeUCR}{RGB}{109,192,103}
\definecolor{yelloUCR}{RGB}{255,224,106}

%%%%%%%%%%%%%%%%%%%%%%%%%%%
%% 3. Theorems and suchlike
%%%%%%%%%%%%%%%%%%%%%%%%%%%

\ifx\nlang\undefined

\theoremstyle{plain}
\ifx \nextra \undefined
\newtheorem{Th}{Teorema}[section]      %%% Theorem 1.1.1
\newtheorem{Tmon}[Th]{Teoremón}
\newtheorem{Prop}[Th]{Proposición}     %%% Proposition 1.1.2
\newtheorem{Lem}[Th]{Lema}             %%% Lemma 1.1.3
\newtheorem{Cor}[Th]{Corolario}        %%% Corollary 1.1.4
\else
\newtheorem{Th}{Teorema}               %%% Theorem 1.1.1
\newtheorem{Tmon}{Teoremón}
\newtheorem{Prop}{Proposición}         %%% Proposition 1.1.2
\newtheorem{Lem}{Lema}                 %%% Lemma 3
\newtheorem{Cor}{Corolario}            %%% Corollary 4
\fi
\newtheorem*{nonum-Th}{Teorema}        %%% No-numbered Theorem
\newtheorem*{nonum-Prop}{Proposición}  %%% No-numbered Proposition
\newtheorem*{nonum-Cor}{Corolario}     %%% No-numbered Corollary

\theoremstyle{definition}
\ifx \nextra \undefined
\newtheorem{Def}[Th]{Definición}       %%% Definition 1.1.5
\newtheorem{Ex}[Th]{Ejemplo}           %%% Example 1.1.6
\newtheorem{Ej}[Th]{Ejercicio}         %%% Ejercicio 1.1.7
\else
\newtheorem{Def}{Definición}           %%% Definition 5
\newtheorem{Ex}{Ejemplo}               %%% Example 6
\newtheorem{Ej}{Ejercicio}             %%% Ejercicio 7
\fi
\newtheorem{Hec}[Th]{Hecho}            %%% Hecho 1.1.8
\newtheorem*{nonum-Def}{Definición}    %%% No number Definition
\newtheorem*{nonum-Ex}{Ejemplo}        %%% No number Example
\newtheorem*{nonum-Ej}{Ejercicio}      %%% No number Ejercicio
\newtheorem*{nonum-Hec}{Hecho}         %%% No number Fact


\theoremstyle{remark}
\newtheorem{Rmk}[Th]{Observación}      %%%Remark 1.1.9
\newtheorem*{nonum-Rmk}{Observación}         %%% No number Fact
\newtheorem*{Notn}{Notaci\'on}        %% Notaciones
\newtheorem*{Warn}{Advertencia}       %% Advertencias
\newtheorem*{Qn}{Pregunta}            %% Pregunta

\else

\theoremstyle{plain}
\ifx \nextra \undefined
\newtheorem{Th}{Theorem}[section]      %%% Theorem 1.1.1
\newtheorem{Tmon}[Th]{Teoremón}
\newtheorem{Prop}[Th]{Proposition}     %%% Proposition 1.1.2
\newtheorem{Lem}[Th]{Lemma}             %%% Lemma 1.1.3
\newtheorem{Cor}[Th]{Corollary}        %%% Corollary 1.1.4
\else
\newtheorem{Th}{Theorem}               %%% Theorem 1.1.1
\newtheorem{Tmon}{Teoremón}
\newtheorem{Prop}{Proposition}         %%% Proposition 1.1.2
\newtheorem{Lem}{Lemma}                 %%% Lemma 3
\newtheorem{Cor}{Corollary}            %%% Corollary 4
\fi
\newtheorem*{nonum-Th}{Theorem}        %%% No-numbered Theorem
\newtheorem*{nonum-Prop}{Proposition}  %%% No-numbered Proposition
\newtheorem*{nonum-Cor}{Corollary}     %%% No-numbered Corollary

\theoremstyle{definition}
\ifx \nextra \undefined
\newtheorem{Def}[Th]{Definition}       %%% Definition 1.1.5
\newtheorem{Ex}[Th]{Example}           %%% Example 1.1.6
\newtheorem{Ej}[Th]{Exercise}         %%% Exercise 1.1.7
\else
\newtheorem{Def}{Definition}           %%% Definition 5
\newtheorem{Ex}{Example}               %%% Example 6
\newtheorem{Ej}{Exercise}             %%% Exercise 7
\fi
\newtheorem{Hec}[Th]{Fact}            %%% Fact 1.1.8
\newtheorem*{nonum-Def}{Definition}    %%% No number Definition
\newtheorem*{nonum-Ex}{Example}        %%% No number Example
\newtheorem*{nonum-Ej}{Exercise}      %%% No number Exercise
\newtheorem*{nonum-Hec}{Fact}         %%% No number Fact


\theoremstyle{remark}
\newtheorem{Rmk}[Th]{Remark}      %%%Remark 1.1.9
\newtheorem*{nonum-Rmk}{Remark}         %%% No number Fact
\newtheorem*{Notn}{Notation}        %% Notaciones
\newtheorem*{Warn}{Warning}       %% Warnings
\newtheorem*{Qn}{Question}            %% Question

\fi 

\numberwithin{equation}{section}

\setlength{\parindent}{3ex}

% \renewcommand{\labelitemi}{--}
% \renewcommand{\labelitemii}{$\circ$}
% \renewcommand{\labelenumi}{(\roman{*})}

%\let\stdsection\section
%\renewcommand\section{\newpage\stdsection}

\newcommand\qedsym{\hfill\ensuremath{\square}}
% Strike through
\newcommand{\st}{\sout}

%%%%%%%%% === My T Color Box === %%%%%%%%%%%%%%

\ifx\nlang\undefined
\ifx \darktheme\undefined
\newtcolorbox{ptcb}{
colframe = black,
colback = white,
breakable,
enhanced
}
\newtcolorbox{ptcbp}{
colframe = black,
colback = white,
coltitle = black,
colbacktitle = black!40,
title = Prueba,
breakable,
enhanced
}
\newtcolorbox{ptcbr}{
colframe = blue,
colback = white,
coltitle = blue,
colbacktitle = blue!40,
title = Respuesta,
breakable,
enhanced
}
\else
\newtcolorbox{ptcb}{
colframe = white,
colback = black,
colupper = white,
breakable,
enhanced
}
\newtcolorbox{ptcbp}{
colframe = white,
colback = dankBlue,
colupper = white,
coltitle = white,
colbacktitle = black,
title = Prueba,
breakable,
enhanced
}
\newtcolorbox{ptcbr}{
colframe = white,
colback = white,
coltitle = blue,
colbacktitle = blue!40,
title = Respuesta,
breakable,
enhanced
}
\fi

\else
\ifx \darktheme\undefined
\newtcolorbox{ptcb}{
colframe = black,
colback = white,
breakable,
enhanced
}
\newtcolorbox{ptcbp}{
colframe = black,
colback = white,
coltitle = black,
colbacktitle = black!40,
title = Proof,
breakable,
enhanced
}
\newtcolorbox{ptcbr}{
colframe = blue,
colback = white,
coltitle = blue,
colbacktitle = blue!40,
title = Answer,
breakable,
enhanced
}
\else
\newtcolorbox{ptcb}{
colframe = white,
colback = black,
colupper = white,
breakable,
enhanced
}
\newtcolorbox{ptcbp}{
colframe = white,
colback = black,
colupper = white,
coltitle = white,
colbacktitle = black!85,
title = Proof,
breakable,
enhanced
}
\newtcolorbox{ptcbr}{
colframe = white,
colback = white,
coltitle = black,
colbacktitle = black!85,
title = Answer,
breakable,
enhanced
}
\fi
\fi


%%%%%%%%% === Listings === %%%%%%%%%%%%%%
\lstset{basicstyle=\ttfamily,breaklines=true}

\lstset{language=Matlab,%
    %basicstyle=\color{red},
    breaklines=true,%
    morekeywords={matlab2tikz},
    keywordstyle=\color{blue},%
    morekeywords=[2]{1}, keywordstyle=[2]{\color{black}},
    identifierstyle=\color{black},%
    stringstyle=\color{MATLABlila},
    commentstyle=\color{MATLABgreen},%
    showstringspaces=false,%without this there will be a symbol in the places where there is a space
    numbers=left,%
    numberstyle={\tiny \color{black}},% size of the numbers
    numbersep=9pt, % this defines how far the numbers are from the text
   % emph=[1]{for,end,break,function,if,elseif,else},emphstyle=[1]\color{blue}, %some words to emphasise
    %emph=[2]{word1,word2}, emphstyle=[2]{style},
}

%%%%%%%%%%%%%%%%%%%%%%%%%%
%% 4. Simple abbreviations
%%%%%%%%%%%%%%%%%%%%%%%%%%

%%% Operator names:

\DeclareMathOperator{\area}{area}
\DeclareMathOperator{\card}{card}
\DeclareMathOperator{\ccl}{ccl}
\DeclareMathOperator{\ch}{ch}
\DeclareMathOperator{\cl}{cl}
\DeclareMathOperator{\coker}{coker}
\DeclareMathOperator{\Conv}{Conv}   %%Convex hull
\DeclareMathOperator{\cosec}{cosec}
\DeclareMathOperator{\cosech}{cosech}
\DeclareMathOperator{\covol}{covol}
\DeclareDocumentCommand\curl{}{\operatorname{curl}} 
\DeclareMathOperator{\diag}{diag}
\DeclareMathOperator{\diam}{diam}
\DeclareMathOperator{\Diff}{Diff}
\DeclareDocumentCommand\div{}{\operatorname{div}} 
\DeclareMathOperator{\energy}{energy}
\DeclareMathOperator{\erfc}{erfc}
\DeclareMathOperator{\Ext}{Ext}
\DeclareMathOperator{\fst}{fst}
\DeclareMathOperator{\Fit}{Fit}
\DeclareMathOperator{\gr}{gr}
\DeclareMathOperator{\hcf}{hcf}
\DeclareMathOperator{\Hilb}{Hilb} %Hilbert scheme
\DeclareMathOperator{\id}{id}
\DeclareMathOperator{\Ind}{Ind}
\DeclareMathOperator{\Int}{Int}
\DeclareMathOperator{\Isom}{Isom}
\DeclareMathOperator{\lcm}{lcm}
\DeclareMathOperator{\length}{length}
\DeclareMathOperator{\Lie}{Lie}
\DeclareMathOperator{\like}{like}
\DeclareMathOperator{\Lk}{Lk}
\DeclareMathOperator{\Maps}{Maps}
\DeclareMathOperator{\mcd}{mcd}
\DeclareMathOperator{\mcm}{mcm}
\DeclareMathOperator{\Min}{Min}
\DeclareMathOperator{\orb}{orb}
\DeclareMathOperator{\ord}{ord}
\DeclareMathOperator{\otp}{otp}
\DeclareMathOperator{\pr}{pr}       %% proyector
\DeclareMathOperator{\poly}{poly}
\DeclareMathOperator{\rel}{rel}
\DeclareMathOperator{\Rad}{Rad}
\DeclareMathOperator*{\res}{res}
\DeclareMathOperator{\Ric}{Ric}
\DeclareMathOperator{\rk}{rk}
\DeclareMathOperator{\Rees}{Rees}
\DeclareMathOperator{\Root}{Root}
\DeclareMathOperator{\rot}{rot}         %% rotacional
\DeclareMathOperator{\spn}{span}
\DeclareMathOperator{\St}{St}
\DeclareMathOperator{\supp}{supp}
\DeclareMathOperator{\Syl}{Syl}
\DeclareMathOperator{\Sym}{Sym}
\DeclareMathOperator{\tot}{tot}     %total Chern class
\DeclareMathOperator{\vol}{vol}

% not-math
\newcommand{\bolds}[1]{{\bfseries #1}}
\newcommand{\cat}[1]{\mathsf{#1}}
\newcommand{\ph}{\,\cdot\,}
\newcommand{\term}[1]{\un{#1}\index{#1}}
\newcommand{\phantomeq}{\hphantom{{}={}}}
\newcommand{\ttt}{\texttt}
\newcommand{\red}[1]{\textcolor{red}{#1}}
\newcommand{\prp}[1]{\textcolor{purple}{#1}}
\newcommand{\blu}[1]{\textcolor{azulUCR}{#1}}
\newcommand{\green}[1]{\textcolor{verdeUCR}{#1}}
\newcommand{\yelo}[1]{\textcolor{yelloUCR}{#1}}
\newcommand{\cele}[1]{\textcolor{celesUCR}{#1}}

%functions
\DeclareMathOperator{\sgn}{sgn}
\newcommand*{\Cdot}{{\raisebox{-0.25ex}{\scalebox{1.5}{$\cdot$}}}}      %% cdot más grande
\newcommand{\ind}{\mathbf{1}}       %%%indicator function
\newcommand{\mm}{\mathfrak{m}}      %%%metric


% Greek letters:

\newcommand{\al}{\alpha}                %% short for  \alpha
\newcommand{\bt}{\beta}                 %% short for  \beta
\newcommand{\Dl}{\Delta}                %% short for  \Delta
\newcommand{\dl}{\delta}                %% short for  \delta
\newcommand{\eps}{\varepsilon}          %% short for  \varepsilon
\newcommand{\Ga}{\Gamma}                %% short for  \Gamma
\newcommand{\ga}{\gamma}                %% short for  \gamma
\newcommand{\kp}{\kappa}                %% short for  \kappa
\newcommand{\La}{\Lambda}               %% short for  \Lambda
\newcommand{\la}{\lambda}               %% short for  \lambda
\newcommand{\Om}{\Omega}                %% short for  \Omega
\newcommand{\om}{\omega}                %% short for  \omega
\newcommand{\Sg}{\Sigma}                %% short for  \Sigma
\newcommand{\sg}{\sigma}                %% short for  \sigma
\newcommand{\Te}{\Theta}                %% short for  \Theta
\newcommand{\te}{\theta}                %% short for  \theta
\newcommand{\ups}{\upsilon}             %% short for  \upsilon
\newcommand{\vf}{\varphi}               %% short for  \varphi
\newcommand{\ze}{\zeta}                 %% short for  \zeta
\newcommand{\vsg}{\varsigma}            %% short for  \varsigma
\newcommand{\vte}{\vartheta}            %% short for  \vartheta

%Boldface letters

\newcommand{\bA}{\mathbb{A}}        %% antisimetrizador
\newcommand{\bB}{\mathbb{B}}        %% bola unitaria
\newcommand{\bC}{\mathbb{C}}    %%% números complejos
\newcommand{\bCP}{\mathbb{CP}}  %%% espacio proyectivo complejo
\newcommand{\bD}{\mathbb{D}}        %% Poincaré disk
\newcommand{\bE}{\mathbb{E}}
\newcommand{\bF}{\mathbb{F}}        %% un cuerpo
\newcommand{\bH}{\mathbb{H}}        %% cuaterniones
\newcommand{\bI}{\mathbb{I}}        %% ideal de zeros
\newcommand{\bK}{\mathbb{K}}            %% ein korper
\newcommand{\bL}{\mathbb{L}}            %% cotangent line bundle
\newcommand{\bN}{\mathbb{N}}    %%% números naturales
\newcommand{\bP}{\mathbb{P}}        %% números enteros positivos
\newcommand{\bQ}{\mathbb{Q}}    %%% números racionales
\newcommand{\bR}{\mathbb{R}}    %%% números reales
\newcommand{\bRP}{\mathbb{RP}}  %%% espacio proyectivo real
\newcommand{\bS}{\mathbb{S}}    %%% esfera
\newcommand{\bT}{\mathbb{T}}        %% círculo o toro
\newcommand{\bV}{\mathbb{V}}        %% lugar geométrico de ceros
\newcommand{\bZ}{\mathbb{Z}}    %%% números enteros

%Script letters:

\newcommand{\cA}{\mathcal{A}}           %% formas diferenciales
\newcommand{\cB}{\mathcal{B}}           %% una base vectorial
\newcommand{\cC}{\mathcal{C}}           %% otra base vectorial
\newcommand{\cD}{\mathcal{D}}           %% funciones de prueba
\newcommand{\cE}{\mathcal{E}}           %% un modulo proyectivo
\newcommand{\cF}{\mathcal{F}}           %% espacio de Fock
\newcommand{\cG}{\mathcal{G}}           %% funtor de Gelfand
\newcommand{\cH}{\mathcal{H}}           %% espacio de Hilbert
\newcommand{\cI}{\mathcal{I}}           %% un funtor de inclusion
\newcommand{\cJ}{\mathcal{J}}           %% otro funtor
\newcommand{\cK}{\mathcal{K}}           %% otro espacio de Hilbert
\newcommand{\cL}{\mathcal{L}}           %% operadores lineales
\newcommand{\cM}{\mathcal{M}}           %% multiplicadores
\newcommand{\cN}{\mathcal{N}}           %% funciones nulas
\newcommand{\cO}{\mathcal{O}}           %% funciones de crec-to lento
\newcommand{\cP}{\mathcal{P}}           %% una particion
\newcommand{\cR}{\mathcal{R}}           %% funciones representativas
\newcommand{\cQ}{\mathcal{Q}}           %% otra particion
\newcommand{\cS}{\mathcal{S}}           %% funciones de Schwartz
\newcommand{\cT}{\mathcal{T}}           %% una topologia
\newcommand{\cU}{\mathcal{U}}           %% cubrimiento abierto
\newcommand{\cV}{\mathcal{V}}           %% vecindarioas
\newcommand{\cW}{\mathcal{W}}           %% grupo de Weyl
\newcommand{\cZ}{\mathcal{Z}}           %% topología de Zariski

%%% Fraktur letters:

\newcommand{\gA}{\mathfrak{A}}      %% un atlas
\newcommand{\g}{\mathfrak{g}}       %% un álgebra de Lie
\newcommand{\gB}{\mathfrak{B}}      %% otro atlas
\newcommand{\ggl}{\mathfrak{gl}}    %% álg de Lie general lineal
\newcommand{\gS}{\mathfrak{S}}      %% grupo simétrico, categoría de esquemas
\newcommand{\gsl}{\mathfrak{sl}}    %% álg de Lie especial lineal
\newcommand{\gso}{\mathfrak{so}}    %% álg de Lie especial ortogonal
\newcommand{\gsu}{\mathfrak{su}}    %% álg de Lie especial unitaria
\newcommand{\gX}{\mathfrak{X}}      %% campos vectoriales

%%% Roman letters:

\newcommand{\dR}{\mathrm{dR}}       %% cohomología de de Rham
\newcommand{\rGL}{\mathrm{GL}}      %% grupo general lineal
\newcommand{\rO}{\mathrm{O}}        %% grupo ortogonal
\newcommand{\rSL}{\mathrm{SL}}      %% grupo especial lineal
\newcommand{\rSO}{\mathrm{SO}}      %% grupo ortogonal especial
\newcommand{\rSp}{\mathrm{Sp}}      %% grupo simpléctico
\newcommand{\rSU}{\mathrm{SU}}      %% grupo unitario especial
\newcommand{\rU}{\mathrm{U}}        %% grupo unitario
\newcommand{\rUH}{\mathrm{UH}}      %% cuaterniones unitarias
\newcommand{\rT}{\mathrm{T}}        %% grupo triangular

% Sanserif letters:

\newcommand{\sA}{\mathsf{A}}            %% algebras de Lie A_n
\newcommand{\sB}{\mathsf{B}}            %% grupo como categoria
\newcommand{\sC}{\mathsf{C}}            %% una categoria
\newcommand{\sD}{\mathsf{D}}            %% otra categoria
\newcommand{\sE}{\mathsf{E}}            %% otra categoria mas
\newcommand{\sF}{\mathsf{F}}            %% algebra de Lie F_4
\newcommand{\sG}{\mathsf{G}}            %% algebra de Lie G_2
\newcommand{\sJ}{\mathsf{J}}            %% un poset
\newcommand{\sK}{\mathsf{K}}            %% un poset
\newcommand{\sL}{\mathcal{L}}           %% derivada de Lie
\newcommand{\sN}{\mathsf{N}}            %% categoría con objetos \bN
\newcommand{\sT}{\mathsf{T}}            %% transpuesta
\newcommand{\Sch}{\cat{Sch}}            %% categoría de esquemas
\newcommand{\Curves}{\cat{Curves}}      %% categoría de curvas
\newcommand{\AlgSp}{\cat{AlgSp}}        %% categoría de espacios algebraicos

%%% Boldface letters:

\bmdefine{\CC}{C}                       %% C negrilla
\bmdefine{\cc}{c}
%\bmdefine{\dd}{d}                       %% d negrilla
\bmdefine{\ee}{e}                       %% vector e
\bmdefine{\eeps}{\varepsilon}           %% basic form \eps
\bmdefine{\FF}{F}                       %% vector F
\bmdefine{\ff}{f}                       %% vector f
\bmdefine{\ii}{i}                       %% cuaternion i
\bmdefine{\jj}{j}                       %% cuaternion j
\bmdefine{\kk}{k}                       %% cuaternion k
\bmdefine{\lla}{\lambda}                %% sucesion \la
\bmdefine{\mmu}{\mu}                    %% sucesion \mu
\bmdefine{\pp}{p}                       %% vector p
\bmdefine{\pt}{pt.}                     %% clase de un punto
\bmdefine{\qq}{q}                       %% vector q
\bmdefine{\rr}{r}                       %% vector r
\bmdefine{\ssg}{\sigma}                 %% vector \sg
%\bmdefine{\sss}{s}
%\bmdefine{\ttt}{t}
\bmdefine{\VV}{V}                       %% V negrilla
\bmdefine{\xx}{x}                       %% sucesion x
\bmdefine{\xxi}{\xi}                    %% vector \xi
\bmdefine{\yy}{y}                       %% sucesion y
\bmdefine{\zz}{z}                       %% sucesion z

% Matrix groups
\DeclareMathOperator{\GL}{GL}   %%% grupo general lineal
\DeclareMathOperator{\Or}{O}    %%% grupo ortogonal
\DeclareMathOperator{\PGL}{PGL} %%% grupo proyectivo lineal
\DeclareMathOperator{\PSL}{PSL} %%% grupo proyectivo lineal especial
\DeclareMathOperator{\PSO}{PSO} %%% grupo proyectivo ortogonal
\DeclareMathOperator{\PSU}{PSU} %%% grupo proyectivo unitario
\DeclareMathOperator{\SL}{SL}   %%% grupo especial lineal
\DeclareMathOperator{\SO}{SO}   %%% grupo especial ortogonal
\DeclareMathOperator{\SU}{SU}   %%% grupo especial unitario

% Numericc
\newcommand{\argmin}{\text{argm\'in}}
\DeclareMathOperator{\dof}{dof}

%% Brackets
\newcommand{\conj}[1]{\left\lbrace#1\right\rbrace}
\newcommand{\bonj}[1]{\left\lbrack#1\right\rbrack}
\newcommand{\obonj}[1]{\left\rbrack#1\right\lbrack}
\newcommand{\rbonj}[1]{\left\rbrack#1\right\rbrack}
\newcommand{\lbonj}[1]{\left\lbrack#1\right\lbrack}
\newcommand{\snm}[1]{\|#1\|}           %small norma
\newcommand{\nm}[1]{\left\|#1\right\|} %norma pegadita
\newcommand{\pnm}[1]{\biggl|\biggl|#1\biggr|\biggr|}
\let\oldvec=\vec
\renewcommand{\vec}[1]{\mathbf{#1}}
\newcommand\quot[2]{
        \mathchoice
            {% \displaystyle
                \text{\raise1ex\hbox{$#1$}\Big/\lower1ex\hbox{$#2$}}%
            }
            {% \textstyle
                {^{ #1}/_{ #2}}
            }
            {% \scriptstyle
                {^{ #1}/_{ #2}}
            }
            {% \scriptscriptstyle
                {^{ #1}/_{ #2}}
            }
    }
%\newcommand*\quot[2]{{^{\textstyle #1}\big/_{\textstyle #2}}}
\newcommand*\squot[2]{{^{ #1}/_{ #2}}}%%%small quotient
\newcommand{\multinom}[2]{\ensuremath{\left(\kern-.3em\left(\genfrac{}{}{0pt}{}{#1}{#2}\right)\kern-.3em\right)}}

% Probability
\DeclareMathOperator{\Bernoulli}{Bernoulli}
\DeclareMathOperator{\betaD}{beta}
\DeclareMathOperator{\bias}{bias}
\DeclareMathOperator{\binomial}{binomial}
\DeclareMathOperator{\corr}{corr}
\DeclareMathOperator{\cov}{cov}
\DeclareMathOperator{\gammaD}{gamma}
\DeclareMathOperator{\mse}{mse}
\DeclareMathOperator{\multinomial}{multinomial}
\DeclareMathOperator{\Poisson}{Poisson}
\DeclareMathOperator{\Var}{Var}     %%%variance
\DeclareMathOperator{\Cov}{Cov}     %%%Covariance
\renewcommand{\mid}{\;\ifnum\currentgrouptype=16 \middle\fi|\;}

% Combinatorics
\DeclareMathOperator{\ins}{ins}   % insertion tableaux
\DeclareMathOperator{\inv}{inv}   % number of inversions
\DeclareMathOperator{\maj}{maj}   % major index
\DeclareMathOperator{\asc}{asc}   % ascents
\DeclareMathOperator{\rw}{rw}     % reading word
\DeclareMathOperator{\rev}{rev}     % reading word
\DeclareMathOperator{\rect}{rect} % rectification of young tableau
\DeclareMathOperator{\sh}{sh}     % shape of young tableau
\DeclareMathOperator{\std}{std}   % standarization
\DeclareMathOperator{\Fl}{\mathcal{F}\ell}       %% conjunto de Flags
\DeclareMathOperator{\Frob}{Frob} % Frobenius map

% Algebra
\DeclareMathOperator{\Ad}{Ad}       %% acción adjunta
\DeclareMathOperator{\adj}{adj}
\DeclareMathOperator{\Ann}{Ann}     %% aniquilador o anulador de módulos
\DeclareMathOperator{\Ass}{Ass}     %% ideales asociados
\DeclareMathOperator{\Aut}{Aut}
\DeclareMathOperator{\Bl}{\mathcal{B}\!\ell}       %% blowup de un espacio
\DeclareMathOperator{\Char}{char}
\DeclareMathOperator{\codim}{codim}
\DeclareMathOperator{\DefB}{Def}     %%% Deformations bundle
%\renewcommand{\Def}{\operatorname{Def}}
\DeclareMathOperator{\disc}{disc}
\DeclareMathOperator{\dom}{dom}
\DeclareMathOperator{\End}{End}     %%%space of endomorphisms
\DeclareMathOperator{\Fix}{Fix}
\DeclareMathOperator{\Frac}{Frac}
\DeclareMathOperator{\Gal}{Gal}
\DeclareMathOperator{\gen}{gen}     %%%set generated by...
\DeclareMathOperator{\Gr}{Gr}       %%%Grassmannian
\DeclareMathOperator{\Hom}{Hom}
\DeclareMathOperator{\Hurw}{Hurw}
\DeclareMathOperator{\image}{image}
\DeclareMathOperator{\Mor}{Mor}
\DeclareMathOperator{\Nil}{Nil}
\DeclareMathOperator{\Ob}{Ob}       %%% Obstructions bundle
\DeclareMathOperator{\Orb}{Orb}
\DeclareMathOperator{\Pf}{Pf}       %%% Pfaffiano de una matrix
\DeclareMathOperator{\Pic}{Pic}     %%% grupo de Picard 
\DeclareMathOperator{\Quot}{Quot}
\DeclareMathOperator{\Spec}{Spec}
\DeclareMathOperator{\Stab}{Stab}
\DeclareMathOperator{\Taut}{Taut}

%Geometry/Topology
\DeclareMathOperator{\glu}{gl}          %% Moduli space gluing morphism
\DeclareMathOperator{\ps}{\mathrm{ps}}       %%% pseudostable superscript for lambdaclass
\newcommand{\vir}{\mathrm{vir}}         %% fibrado virtual

% Analysis
\DeclareMathOperator*{\esssup}{ess\hspace{0.5mm}sup}
\DeclareMathOperator*{\essinf}{ess\hspace{0.5mm}inf}
%\DeclareMathOperator{\Int}{Int}     %%%interior vacilon funcional

\newcommand{\loc}{\text{loc}}
\newcommand{\LB}{\cL_\cB}           %%%bounded linear operator

% Logic
\newcommand{\cleq}{\preccurlyeq}
\newcommand{\cgeq}{\succcurlyeq}

% Others
\renewcommand{\ev}{\operatorname{ev}}     %%%evalutation previously expectation value physics package
\newcommand{\bigcupdot}{\charfusion[\mathop]{\bigcup}{\Cdot}} %%JCVDG
%\renewcommand{\bigcupdot}{\charfusion[\mathop]{\bigcup}{\Cdot}}
\newcommand{\cupdot}{\charfusion[\mathbin]{\cup}{\Cdot}}
\newcommand{\exterior}{\mathchoice{{\textstyle\bigwedge}}{{\bigwedge}}{{\textstyle\wedge}}{{\scriptstyle\wedge}}}
\newcommand{\hol}{\mathfrak{hol}}
\newcommand{\Id}{\mathrm{Id}}
\newcommand{\lie}[1]{\mathfrak{#1}}
\newcommand{\qeq}{\mathrel{``{=}"}}
\newcommand{\wsto}{\stackrel{\mathrm{w}^*}{\to}}
\newcommand{\wt}{\mathrm{wt}}

%\let\Im\relax
%\let\Re\relax

%%% Shorter symbol names:

\newcommand{\bull}{{\scriptstyle\bullet}}  %% vertice en figuras
\newcommand{\del}{\partial}             %% short for  \partial
\newcommand{\downto}{\downarrow}        %% limite a la derecha
\newcommand{\dsp}{\displaystyle}        %% despliegue en texto
\renewcommand{\geq}{\geqslant}          %% mayor o igual (variante)
\newcommand{\hookto}{\hookrightarrow}     %% inclusion arrow
\newcommand{\isom}{\simeq}              %% isomorfismo
\renewcommand{\l}{\ell}                   %% ele cursiva
\renewcommand{\leq}{\leqslant}          %% menor o igual (variante)
\newcommand{\less}{\setminus}           %% set difference
\newcommand{\otto}{\leftrightarrow}     %% bijection
\newcommand{\ox}{\otimes}               %% producto tensorial
\newcommand{\rt}{\triangleleft}         %% un orden parcial
\newcommand{\rteq}{\trianglelefteq}     %% normal subgroup
\newcommand{\up}{{\mathord{\uparrow}}}  %% espinor `up'
\newcommand{\upto}{\uparrow}            %% left hand limit
\newcommand{\w}{\wedge}                 %% producto exterior
\newcommand{\wto}{\rightharpoonup}      %% convergencia debil
\newcommand{\x}{\times}                 %% producto vectorial
\renewcommand{\.}{\Cdot}                %% producto escalar
\renewcommand{\:}{\mathbin{:}}          %% colon in  f: A -> B
\newcommand{\into}{\rightarrowtail}     %% injection arrow
\newcommand{\lr}{\dashv}                %% adjunction
\newcommand{\lt}{\triangleright}        %% a left action
\newcommand{\lteq}{\trianglerighteq}    %% normal supergroup
\newcommand{\nb}{\nabla}                %% homomorfismo de suma
\newcommand{\nisom}{\not\simeq}         %% negacion de isomorfismo
%\newcommand{\oast}{\circledast}         %% variante de * (ya existe en stmaryrd)
\newcommand{\onto}{\twoheadrightarrow}  %% surjection arrow
\newcommand{\opp}{\circ}                %% objeto opuesto
\newcommand{\ottto}{\longleftrightarrow} %% bijection in display
\newcommand{\pullb}{\lrcorner}          %% simbolo de pullback
\newcommand{\pushf}{\ulcorner}          %% simbolo de pushout
\newcommand{\rx}{\rtimes}               %% producto semidirecto
\newcommand{\To}{\Rightarrow}           %% entre funtores
\newcommand{\tofro}{\rightleftarrows}   %% pair of opposed maps
\newcommand{\toto}{\rightrightarrows}   %% pair of parallel maps

\renewcommand{\2}{\flat}                  %% marcador de sucesiones
\newcommand{\3}{\sharp}                 %% marcador de sucesiones
\newcommand{\4}{\natural}               %% marcador de morfismos
% \newcommand{\5}{\diamond}               %% for roots of trees
% \newcommand{\7}{\dagger}                %% adjunto de operador
\newcommand{\8}{\bullet}                %% anonymous degree

%%% Useful abbreviations:

\newcommand{\Coo}{\cC^\infty}         %% funciones suaves
\newcommand{\ctr}{\mathbin{\lrcorner\,}} %% contraction symbol
\newcommand{\nbf}{{\vec\nabla}}     %% short for  \vec\nabla

\newcommand{\as}{\quad\text{cuando}\enspace} %% `cuando' en límites
\newcommand{\bCoo}{{\bC_\infty}}    %% esfera de Riemann
% \newcommand{\bRoo}{{\bR_\infty}}    %% círculo real extendido

%%% Repeated relations:

\newcommand{\cupycup}{\cup\cdots\cup} %% unión repetida
\newcommand{\capycap}{\cap\cdots\cap} %% intersección repetida
\newcommand{\sys}{\subset\cdots\subset}%% subconjunto propio repetido
\newcommand{\subysub}{\subseteq\cdots\subseteq} %%subconjunto repetido
\newcommand{\oxyox}{\otimes\cdots\otimes} %% prod tensorial repetido
\newcommand{\wyw}{\wedge\cdots\wedge} %% producto exterior repetido
\newcommand{\opyop}{\oplus\cdots\oplus} %% suma directa repetida
\newcommand{\xyx}{\times\cdots\times} %% producto directo repetido

%%% Arrows with riders:

\newcommand{\longto}{\mathop{\longrightarrow}\limits}

%%% Small fractions in displays:

\newcommand{\half}{{\mathchoice{\nhalf}{\thalf}{\shalf}{\shalf}}} %%display text script script^2
\newcommand{\happi}{{\tfrac{\pi}{2}}} %% small fraction  \pi/2
\newcommand{\quarter}{\tfrac{1}{4}} %% small fraction  1/4
\newcommand{\nhalf}{\frac{1}{2}}
\newcommand{\shalf}{{\scriptstyle\frac{1}{2}}} %% tiny fraction 1/2
\newcommand{\thalf}{{\tfrac{1}{2}}} %% small fraction  1/2
\renewcommand{\third}{\tfrac{1}{3}}   %% small fraction  1/3 %Hay que renew porque mathabx toma second y third como x'' y x''' por ejemplo

\newcommand{\ihalf}{{\tfrac{i}{2}}} %% small fraction  i/2

%%%%%%%%%%%%%%%%%%%%%%%%%%%%%
%% 5. Commands with arguments
%%%%%%%%%%%%%%%%%%%%%%%%%%%%%

%%% Accent-like commands, abbreviated:

\newcommand{\ov}{\overline}        %% short for  \overline
\newcommand{\un}{\underline}       %% short for  \underline
\newcommand{\wh}{\widehat}          %% short for  \widehat

%%% Separate words in displays:

\newcommand{\word}[1]{\quad\text{#1}\quad} %% texto intercalado

%%% Webpage locator:

\newcommand{\zelda}[1]{$\langle${\footnotesize\texttt{#1}}$\rangle$}

%% Symbol placement:

\newcommand{\pre}[1]{{}^{#1\!}} %% upper left exponent

%%% Proof-part labels:

\newcommand{\Adiff}[2]{\ensuremath{\Ad\,(\mathrm{#1})\Longleftrightarrow
    (\mathrm{#2})}:\enspace}
\newcommand{\Adimp}[2]{\ensuremath{\Ad\,(\mathrm{#1})\Longrightarrow
    (\mathrm{#2})}:\enspace}
\newcommand{\Adit}[1]{\ensuremath{\Ad\,(\mathrm{#1})}:\enspace}

%%% Enclose one argument with delimiters:

\newcommand{\bool}[1]{\llbracket#1\rrbracket} %% condición booleana
\newcommand{\combo}[1]{\operatorname{co}(#1)} %% convex combo
\newcommand{\lin}[1]{\operatorname{lin}\langle#1\rangle} %% `span'
\newcommand{\set}[1]{\{\,#1\,\}}    %% set notation

\newcommand{\floor}[1]{\lfloor#1\rfloor} %% mayor entero <= x
\newcommand{\Set}[1]{\biggl\{\,#1\,\biggr\}} %% set notation (large)
\newcommand{\roof}[1]{\lceil#1\rceil} %% menor entero >= x
\newcommand{\genr}[1]{\left\langle #1\right\rangle}     %% grupo generado por #1

%%% Asides:

\newcommand{\aside}[1]{$\llbracket$\,#1\,$\rrbracket$} % nota lateral
\ifx \nlang \undefined
\newcommand{\hint}[1]{$\llbracket$\,In\-di\-ca\-ci\'on: #1\,$\rrbracket$}
\else
\newcommand{\hint}[1]{$\llbracket$\,Hint: #1\,$\rrbracket$}
\fi 


%%% Matrices:

\newcommand{\onebytwo}[2]{\begin{pmatrix} %% 1 x 2 matrix
  #1 & #2 \end{pmatrix}}
\newcommand{\onebythree}[3]{\begin{pmatrix} %% 1 x 3 matrix
  #1 & #2 & #3 \end{pmatrix}}
\newcommand{\onebyfour}[4]{\begin{pmatrix} %% 1 x 4 matrix
  #1 & #2 & #3 & #4 \end{pmatrix}}
\newcommand{\twobyone}[2]{\begin{pmatrix} %% 2 x 1 matrix
   #1 \\ #2 \end{pmatrix}}
\newcommand{\twobytwo}[4]{\begin{pmatrix} %% 2 x 2 matrix
   #1 & #2 \\ #3 & #4 \end{pmatrix}}
\newcommand{\twobythree}[6]{\begin{pmatrix} %% 2 x 3 matrix
    #1 & #2 & #3\\ #4 & #5 & #6 \end{pmatrix}}
\newcommand{\threebyone}[3]{\begin{pmatrix} %% 3 x 1 matrix
   #1 \\ #2 \\ #3 \end{pmatrix}}
\newcommand{\threebythree}[9]{\begin{pmatrix} %% 3 x 3 matrix
   #1 & #2 & #3 \\ #4 & #5 & #6 \\ #7 & #8 & #9 \end{pmatrix}}
\newcommand{\fourbyone}[4]{\begin{pmatrix} %% 2 x 1 matrix
   #1 \\ #2 \\ #3 \\ #4 \end{pmatrix}}
%\newcommand{\fourbyfour}[16]{\begin{pmatrix} %% 4 x 4 matrix
%  #1 & #2 & #3 & #4\\ #5 & #6 & #7 & #8 \\ #9 & #10 & #11 & #12 \\ #13 & #14 & #15 & #16 \end{pmatrix}}
\newcommand{\nbyn}[9]{\begin{pmatrix} %% 4 x 4 matrix with prefilled entries
  #1 & #2 & \cdots & #3\\ #4 & #5 & \cdots & #6 \\ \vdots & \vdots & \ddots & \vdots \\ #7 & #8 & \cdots & #9 \end{pmatrix}}

%%%%%%%%%%%%%%%%%%%%%%%%%%%%
%% 6. Hyphenation exceptions
%%%%%%%%%%%%%%%%%%%%%%%%%%%%

\hyphenation{auto-va-lor auto-va-lo-res auto-vec-tor auto-vec-to-res
car-di-na-li-dad ce-rra-da ce-rra-do ce-rra-das ce-rra-dos cons-tan-te
cons-tan-tes cons-truc-ci cons-truir con-ti-nua con-ti-nua-mente
con-ti-nuas con-ti-nui-dad con-ti-nuo con-ti-nuos co-rres-pon-den-cia
co-rres-pon-de co-rres-pon-den co-rres-pon-dien-te
co-rres-pon-dien-tes co-va-rian-te cual-quier cual-quiera
cu-bri-mien-to desa-rro-lla-do desa-rro-llar des-pu dia-go-nal
dia-go-na-les di-fe-ren-cia-ble di-fe-ren-cia-bles di-fe-ren-cial
di-fe-ren-cia-les di-fe-ren-te di-fe-ren-tes dis-cre-ta dis-cre-tas
dis-cre-to dis-cre-tos di-vi-si-bi-li-dad di-vi-si-ble ele-men-tal
ele-men-ta-les ele-men-to ele-men-tos equi-va-len-cia equi-va-lente
equi-va-lentes equi-va-rian-te equi-va-rian-tes eu-cli-dia-na
eu-cli-dia-nas eu-cli-dia-no eu-cli-dia-nos Fi-gu-ra Gal-ois
gal-oi-sia-na ge-ne-rada ge-ne-rado ge-ne-ra-dor ge-ne-ra-do-res
ge-ne-ral ge-ne-ra-les ge-ne-ra-li-dad ge-ne-ra-li-za ge-ne-ra-li-zan
ge-ne-ran ge-ne-rar geo-me-tr geo-me-try Ha-da-mard ho-meo-mor-fis-mo
ho-meo-mor-fo idea-les in-de-pen-dien-te in-de-pen-dien-tes
in-va-rian-cia in-va-rian-te in-va-rian-tes li-ne-a-les
li-ne-al-men-te ma-ne-ra me-dian-te mo-der-no nin-gu-no nues-tra
nues-tro nu-me-ra-ble ope-ra-ci ope-ra-cio-nes ope-ra-dor
ope-ra-do-res or-to-go-nal par-ti-cu-lar pro-ce-di-mien-to pro-duc-to
pro-duc-tos pro-pie-dad pro-pie-da-des pro-po-si-ci re-fe-ren-cia
re-fle-xi-va re-fle-xi-vas re-fle-xi-vo re-fle-xi-vos re-so-lu-ble
res-pec-ti-va-men-te res-pec-ti-vo res-pec-ti-vos res-pec-to
sa-tis-fa-ce sepa-ra-ble sepa-ra-bles si-guien-te si-guien-tes
subes-pa-cio subes-pa-cios te-dra-edro te-tra-edros tri-vial
tri-via-les uti-lidad va-lo-res va-ria-ble va-ria-bles va-rie-dad
va-rie-da-des ve-cin-da-rio ve-cin-da-rios vec-to-rial vec-to-ria-les
vice-versa}


%%% TikZ arrows and such

\pgfarrowsdeclarecombine{twolatex'}{twolatex'}{latex'}{latex'}{latex'}{latex'}
\tikzset{->/.style = {decoration={markings,
                                  mark=at position 1 with {\arrow[scale=2]{latex'}}},
                      postaction={decorate}}}
\tikzset{<-/.style = {decoration={markings,
                                  mark=at position 0 with {\arrowreversed[scale=2]{latex'}}},
                      postaction={decorate}}}
\tikzset{<->/.style = {decoration={markings,
                                   mark=at position 0 with {\arrowreversed[scale=2]{latex'}},
                                   mark=at position 1 with {\arrow[scale=2]{latex'}}},
                       postaction={decorate}}}
\tikzset{->-/.style = {decoration={markings,
                                   mark=at position #1 with {\arrow[scale=2]{latex'}}},
                       postaction={decorate}}}
\tikzset{-<-/.style = {decoration={markings,
                                   mark=at position #1 with {\arrowreversed[scale=2]{latex'}}},
                       postaction={decorate}}}
\tikzset{->>/.style = {decoration={markings,
                                  mark=at position 1 with {\arrow[scale=2]{latex'}}},
                      postaction={decorate}}}
\tikzset{<<-/.style = {decoration={markings,
                                  mark=at position 0 with {\arrowreversed[scale=2]{twolatex'}}},
                      postaction={decorate}}}
\tikzset{<<->>/.style = {decoration={markings,
                                   mark=at position 0 with {\arrowreversed[scale=2]{twolatex'}},
                                   mark=at position 1 with {\arrow[scale=2]{twolatex'}}},
                       postaction={decorate}}}
\tikzset{->>-/.style = {decoration={markings,
                                   mark=at position #1 with {\arrow[scale=2]{twolatex'}}},
                       postaction={decorate}}}
\tikzset{-<<-/.style = {decoration={markings,
                                   mark=at position #1 with {\arrowreversed[scale=2]{twolatex'}}},
                       postaction={decorate}}}

\tikzset{circ/.style = {fill, circle, inner sep = 0, minimum size = 3}}
\tikzset{scirc/.style = {fill, circle, inner sep = 0, minimum size = 1.5}}
\tikzset{mstate/.style={circle, draw, blue, text=black, minimum width=0.7cm}}

\tikzset{eqpic/.style={baseline={([yshift=-.5ex]current bounding box.center)}}}
\tikzset{commutative diagrams/.cd,cdmap/.style={/tikz/column 1/.append style={anchor=base east},/tikz/column 2/.append style={anchor=base west},row sep=tiny}}

\definecolor{mblue}{rgb}{0.2, 0.3, 0.8}
\definecolor{morange}{rgb}{1, 0.5, 0}
\definecolor{mgreen}{rgb}{0.1, 0.4, 0.2}
\definecolor{mred}{rgb}{0.5, 0, 0}

\def\drawcirculararc(#1,#2)(#3,#4)(#5,#6){%
    \pgfmathsetmacro\cA{(#1*#1+#2*#2-#3*#3-#4*#4)/2}%
    \pgfmathsetmacro\cB{(#1*#1+#2*#2-#5*#5-#6*#6)/2}%
    \pgfmathsetmacro\cy{(\cB*(#1-#3)-\cA*(#1-#5))/%
                        ((#2-#6)*(#1-#3)-(#2-#4)*(#1-#5))}%
    \pgfmathsetmacro\cx{(\cA-\cy*(#2-#4))/(#1-#3)}%
    \pgfmathsetmacro\cr{sqrt((#1-\cx)*(#1-\cx)+(#2-\cy)*(#2-\cy))}%
    \pgfmathsetmacro\cA{atan2(#2-\cy,#1-\cx)}%
    \pgfmathsetmacro\cB{atan2(#6-\cy,#5-\cx)}%
    \pgfmathparse{\cB<\cA}%
    \ifnum\pgfmathresult=1
        \pgfmathsetmacro\cB{\cB+360}%
    \fi
    \draw (#1,#2) arc (\cA:\cB:\cr);%
}
\newcommand\getCoord[3]{\newdimen{#1}\newdimen{#2}\pgfextractx{#1}{\pgfpointanchor{#3}{center}}\pgfextracty{#2}{\pgfpointanchor{#3}{center}}}

\newcommand\qedshift{\vspace{-17pt}}
\newcommand\fakeqed{\pushQED{\qed}\qedhere}

\def\Xint#1{\mathchoice
   {\XXint\displaystyle\textstyle{#1}}%
   {\XXint\textstyle\scriptstyle{#1}}%
   {\XXint\scriptstyle\scriptscriptstyle{#1}}%
   {\XXint\scriptscriptstyle\scriptscriptstyle{#1}}%
   \!\int}
\def\XXint#1#2#3{{\setbox0=\hbox{$#1{#2#3}{\int}$}
     \vcenter{\hbox{$#2#3$}}\kern-.5\wd0}}
\def\ddashint{\Xint=}
\def\dashint{\Xint-}

\newcommand\separator{{\centering\rule{2cm}{0.2pt}\vspace{2pt}\par}}

\newenvironment{own}{\color{gray!70!black}}{}

\newcommand\makecenter[1]{\raisebox{-0.5\height}{#1}}

\mathchardef\mdash="2D

\newenvironment{significant}{\begin{center}\begin{minipage}{0.9\textwidth}\centering\em}{\end{minipage}\end{center}}
\DeclareRobustCommand{\rvdots}{%
  \vbox{
    \baselineskip4\p@\lineskiplimit\z@
    \kern-\p@
    \hbox{.}\hbox{.}\hbox{.}
  }}
\DeclareRobustCommand\tph[3]{{\texorpdfstring{#1}{#2}}}
\def\BState{\State\hskip-\ALG@thistlm}

\makeatother 

% The local additions match the working article notes. The shared header
% supplies the theorem environments, ptcb boxes, Sch, Curves, and isom.
\usepackage{mathtools}
\usepackage{tikz-cd}
\usetikzlibrary{babel}
\tikzcdset{
  every diagram/.append style={font=\large,row sep=large,column sep=large},
  labels={font=\normalsize},
  arrows={line width=.6pt}
}

\let\headerst\st
\let\st\relax
\usepackage{soul}
\let\st\headerst
\soulregister\ref7
\counterwithout{section}{chapter}
\counterwithout{figure}{chapter}
\setsecnumdepth{subsection}
\settocdepth{section}
% Short running heads keep long section names from colliding with the title.
\makeoddhead{ruled}{\small MMARGS}{}{\small\rightmark}
\makeevenhead{ruled}{\small MMARGS}{}{\small\rightmark}

\definecolor{noteblue}{HTML}{00C0F3}
\definecolor{notegreen}{HTML}{6DC067}
\definecolor{notepink}{HTML}{F07DAA}
\makeatletter
\newcommand{\roundedhighlight}[1]{%
  \begingroup
  \setuloverlap{2pt}%
  \def\SOUL@ulunderline##1{{%
    \setbox\z@\hbox{##1}%
    \rlap{\SOUL@ulcolor{\tikz[overlay,blend mode=darken]{%
      \fill[rounded corners=2pt]
        (-\SOUL@uloverlap,-\SOUL@uldp)
        rectangle (\dimexpr\wd\z@+\SOUL@uloverlap\relax,\SOUL@ulht);%
    }}}%
    \unhcopy\z@
  }}%
  \hl{#1}%
  \endgroup
}
\makeatother
\DeclareRobustCommand{\hlblue}[1]{{\sethlcolor{noteblue}\roundedhighlight{#1}}}
\DeclareRobustCommand{\hlgreen}[1]{{\sethlcolor{notegreen}\roundedhighlight{#1}}}
\DeclareRobustCommand{\hlpink}[1]{{\sethlcolor{notepink}\roundedhighlight{#1}}}
\newcounter{passage}
\newcommand{\passagemark}{%
  \refstepcounter{passage}\textsuperscript{\normalfont[\thepassage]}%
}
% Native TikZ diagrams copied from Stacks4StupidPeople.tex.
% Required in the parent preamble: tikz with arrows.meta, the color
% noteblue defined as HTML 00C0F3, and \isom defined as \simeq.
% Object and arrow-label fonts, geometry, and line widths match the notes.
% These macros contain only displays; they introduce no cross-reference labels.

\newcommand{\TalkDDDiagram}{%
\[
\begin{tikzpicture}[
    y=.86cm,
    arrows={-{Latex[length=1.6mm]}},
    line width=.6pt,
    every node/.style={inner sep=2pt,font=\large},
    projection/.style={densely dotted,draw=gray!75},
    comparison/.style={draw=noteblue,line width=.8pt}
]
    % The covering and its overlap form the floor.
    \node (Si)  at (0,0)       {$S_i$};
    \node (Sij) at (2.8,1.4)   {$S_{ij}$};
    \node (Sj)  at (7,.6)      {$S_j$};
    \node (S)   at (4.2,-.8)   {$S$};

    % Descent data supplies these local objects, with no global E.
    \node (Ei)   at (0,4.7)    {$E_i$};
    \node (Eiij) at (1.2,5.9)  {$E_i\mid S_{ij}$};
    \node (Ejij) at (4.4,6.3)  {$E_j\mid S_{ij}$};
    \node (Ej)   at (7,5.3)    {$E_j$};

    \begin{scope}[every node/.append style={font=\normalsize}]
    \draw (Sij) -- node[above left] {$p_i$} (Si);
    \draw (Sij) -- node[pos=.72,above] {$p_j$} (Sj);
    \draw (Sij) -- node[pos=.55,left,xshift=-.2em] {$v$} (S);

    % Dotted connectors depict the action of the single functor pi.
    \draw[projection] (Eiij) -- (Sij);
    \draw[projection] (Ejij) -- (Sij);
    \draw[projection] (Ei) -- (Si);
    \draw[projection] (Ej) -- (Sj);

    \draw (Eiij) --
        node[left,xshift=-.25em,yshift=.4em] {$r_i$} (Ei);
    \draw (Ejij) --
        node[pos=.55,above right,xshift=.2em] {$r_j$} (Ej);
    \draw[comparison] (Eiij) --
        node[above,yshift=.4em] {$\alpha_{ij}$}
        node[below,yshift=-.2em] {$\isom$} (Ejij);

    \draw (Si) -- node[below,yshift=-.2em] {$u_i$} (S);
    \draw (Sj) -- node[below,yshift=-.2em] {$u_j$} (S);
    \end{scope}
\end{tikzpicture}
\]%
}

\newcommand{\TalkDDEDiagram}{%
\[
\begin{tikzpicture}[
    y=.86cm,
    arrows={-{Latex[length=1.6mm]}},
    line width=.6pt,
    every node/.style={inner sep=2pt,font=\large},
    projection/.style={densely dotted,draw=gray!75},
    comparison/.style={draw=noteblue,line width=.8pt},
    foreground/.style={preaction={draw=white,line width=3pt,arrows=-}}
]
    % The floor sits directly below the square of E-pullbacks:
    % each base vertex is translated by (0,-4.4) from its lift.
    \node (Si)  at (3,.45)     {$S_i$};
    \node (Sij) at (5,2.1)     {$S_{ij}$};
    \node (Sj)  at (8,.95)     {$S_j$};
    \node (S)   at (6,-.7)     {$S$};

    % The upper face separates the supplied local objects from
    % the pullbacks of a global object realizing the descent datum.
    \node (Ei)   at (0,6)       {$E_i$};
    \node (Eiij) at (2,7.65)    {$E_i\mid S_{ij}$};
    \node (Ejij) at (7,8.15)    {$E_j\mid S_{ij}$};
    \node (Ej)   at (10,7)      {$E_j$};
    \node (C)    at (5,6.5)     {$E\mid S_{ij}$};
    \node (Bi)   at (3,4.85)    {$E\mid S_i$};
    \node (Bj)   at (8,5.35)    {$E\mid S_j$};
    \node (E)    at (6,3.7)     {$E$};

    \begin{scope}[every node/.append style={font=\normalsize}]
    \draw (Sij) -- node[above left] {$p_i$} (Si);
    \draw (Sij) -- node[pos=.72,above] {$p_j$} (Sj);
    \draw (Si) -- node[below,yshift=-.2em] {$u_i$} (S);
    \draw (Sj) -- node[below,yshift=-.2em] {$u_j$} (S);

    % Dotted connectors show the action of pi on objects.
    \draw[projection] (Eiij) -- (Sij);
    \draw[projection] (Ejij) -- (Sij);
    \draw[projection] (C) -- (Sij);
    \draw[projection] (Ei) -- (Si);
    \draw[projection] (Bi) -- (Si);
    \draw[projection] (Ej) -- (Sj);
    \draw[projection] (Bj) -- (Sj);
    \draw[projection,foreground] (E) -- (S);

    % Existing pullback maps and the two newly named maps s_i,s_j.
    \draw[foreground] (Eiij) --
        node[left,xshift=-.25em,yshift=.4em] {$r_i$} (Ei);
    \draw[foreground] (Ejij) --
        node[pos=.55,above right,yshift=.2em] {$r_j$} (Ej);
    \draw[foreground] (C) --
        node[midway,above left,xshift=-.15em,yshift=.15em] {$s_i$} (Bi);
    \draw[foreground] (C) --
        node[midway,above right,xshift=.15em,yshift=.15em] {$s_j$} (Bj);
    \draw[foreground] (Bi) --
        node[pos=.45,below,yshift=-.2em] {$q_i$} (E);
    \draw[foreground] (Bj) --
        node[below,yshift=-.2em] {$q_j$} (E);

    % Every blue arrow is labeled by its map and an isomorphism symbol.
    \draw[foreground,comparison] (Eiij) --
        node[above,yshift=.3em] {$\alpha_{ij}$}
        node[below,yshift=-.2em] {$\isom$} (Ejij);
    \draw[foreground,comparison] (C) --
        node[pos=.3,above,yshift=.3em] {$\isom$}
        node[pos=.55,below,yshift=-.2em] {$\alpha_i|_{S_{ij}}$} (Eiij);
    \draw[foreground,comparison] (C) --
        node[pos=.32,right,xshift=.8em,yshift=-.15em] {$\alpha_j|_{S_{ij}}$}
        node[pos=.6,left,xshift=-.1em] {$\isom$} (Ejij);
    \draw[foreground,comparison] (Bi) --
        node[pos=.4,above,yshift=.25em] {$\alpha_i$}
        node[pos=.55,below,yshift=-.2em] {$\isom$} (Ei);
    \draw[foreground,comparison] (Bj) --
        node[pos=.45,above,yshift=.25em] {$\alpha_j$}
        node[pos=.4,below,yshift=-.2em] {$\isom$} (Ej);
    \end{scope}
\end{tikzpicture}
\]%
}

\newcommand{\TalkISDiagram}{%
\[
\begin{tikzpicture}[
    % Orthographic view from t(2,-3,1), with z pointing upward.
    % Screen coordinates: ((3x+2y)/sqrt(13),(-2x+3y+13z)/sqrt(182)).
    % The common physical scale is .8 cm per projected unit.
    x={(.665640cm,-.118600cm)},
    y={(.443760cm,.177900cm)},
    z={(0cm,.770899cm)},
    arrows={-{Latex[length=1.6mm]}},
    line width=.6pt,
    every node/.style={inner sep=2pt,font=\large},
    projection/.style={densely dotted,draw=gray!75},
    comparison/.style={draw=noteblue,line width=.8pt}
]
    \pgfdeclarelayer{isom-objects}
    \pgfsetlayers{main,isom-objects}
    % Object labels remain legible above intersecting arrows.
    \begin{pgfonlayer}{isom-objects}
    % A parallel scheme floor, centered below each E/E' pair.
    \node[fill=white] (Si)  at (1.25,1.25,-3) {$S_i$};
    \node[fill=white] (S)   at (10.25,1.25,-3) {$S$};
    \node[fill=white] (Sij) at (1.25,10.25,-3) {$S_{ij}$};
    \node[fill=white] (Sj)  at (10.25,10.25,-3) {$S_j$};

    % The requested upper square: A, A+9i, A+9j, A+9i+9j.
    \node[fill=white] (Bi) at (0,0,2) {$E\mid S_i$};
    \node[fill=white] (E)  at (9,0,2) {$E$};
    \node[fill=white] (A)  at (0,9,2) {$E\mid S_{ij}$};
    \node[fill=white] (Bj) at (9,9,2) {$E\mid S_j$};
    % Its primed copy is translated by (2.5,2.5,0) in the same plane.
    \node[fill=white] (Bip) at (2.5,2.5,2) {$E'\mid S_i$};
    \node[fill=white] (Ep)  at (11.5,2.5,2) {$E'$};
    \node[fill=white] (Ap)  at (2.5,11.5,2) {$E'\mid S_{ij}$};
    \node[fill=white] (Bjp) at (11.5,11.5,2) {$E'\mid S_j$};
    \end{pgfonlayer}

    \begin{scope}[every node/.append style={font=\normalsize,fill=white,inner sep=1pt}]
    \draw (Sij) -- node[above left,xshift=-.15em,yshift=.1em] {$p_i$} (Si);
    \draw (Sij) -- node[above,yshift=.25em] {$p_j$} (Sj);
    \draw (Si) -- node[below,yshift=-.25em] {$u_i$} (S);
    \draw (Sj) -- node[below right,xshift=.15em,yshift=-.1em] {$u_j$} (S);

    % Dotted connectors indicate the projection functor on objects.
    \draw[projection] (A) -- (Sij);
    \draw[projection] (Ap) -- (Sij);
    \draw[projection] (Bi) -- (Si);
    \draw[projection] (Bip) -- (Si);
    \draw[projection] (Bj) -- (Sj);
    \draw[projection] (Bjp) -- (Sj);
    \draw[projection] (E) -- (S);
    \draw[projection] (Ep) -- (S);

    % All black pullback maps are straight.
    \draw (Ap) --
        node[pos=.9,below right,xshift=.4em,yshift=-.1em] {$r_i'$} (Bip);
    \draw (Ap) --
        node[pos=.5,above,yshift=.3em] {$r_j'$} (Bjp);
    \draw (Bip) --
        node[pos=.65,above,yshift=.3em] {$q_i'$} (Ep);
    \draw (Bjp) --
        node[pos=.55,below right,xshift=.25em,yshift=-.2em] {$q_j'$} (Ep);
    \draw (A) --
        node[pos=.5,above left,xshift=-.25em,yshift=.2em] {$r_i$} (Bi);
    \draw (A) --
        node[pos=.6,below,yshift=-.3em] {$r_j$} (Bj);
    \draw (Bi) --
        node[pos=.5,below,yshift=-.3em] {$q_i$} (E);
    \draw (Bj) --
        node[pos=.55,above left,xshift=-.25em,yshift=.2em] {$q_j$} (E);

    % The same overlap morphism is depicted twice, as a visible lens.
    \draw[comparison] (A) to[bend left=50]
        node[pos=.5,above,yshift=.3em]
            {$\alpha_i|_{S_{ij}}\;\isom$} (Ap);
    \draw[comparison] (A) to[bend right=50]
        node[pos=.5,above,yshift=.25em]
            {$\alpha_j|_{S_{ij}}\;\isom$} (Ap);
    \node[fill=white,inner sep=1pt] at (1.25,10.25,2) {$=$};

    \draw[comparison] (Bi) --
        node[pos=.5,below,yshift=-.35em] {$\alpha_i\;\isom$} (Bip);
    \draw[comparison] (Bj) --
        node[pos=.5,above,yshift=.8em] {$\alpha_j\;\isom$} (Bjp);
    % The dashed global isomorphism is supplied by the sheaf condition.
    \draw[comparison,dashed] (E) --
        node[pos=.5,below,yshift=-.35em] {$\alpha\;\isom?$} (Ep);
    \end{scope}
\end{tikzpicture}
\]%
}

\hypersetup{
  pdftitle={Stacks: de las familias a los espacios de módulos -- MMARGS, 29 de octubre de 2026},
  pdfauthor={Ignacio Rojas},
  pdfsubject={Notas de presentación basadas en Stacks for everybody, de Barbara Fantechi}
}

\begin{document}
\thispagestyle{plain}

\begin{center}
  {\LARGE\bfseries Stacks: de las familias\\ a los espacios de módulos}\par
  \vspace{.5em}
  {\large MMARGS --- 29 de octubre de 2026}\par
  \vspace{.35em}
  Ignacio Rojas\par
  {\small\itshape Borrador para la presentación, basado en nuestro estudio de
  \emph{Stacks for everybody}, de Barbara Fantechi.}
\end{center}

\begin{ptcb}
\textbf{La idea que queremos conservar.}
Un stack es una categoría de objetos geométricos que
\hlgreen{admiten cambios de base como las familias} y
\hlblue{se pegan como las familias}. Un stack algebraico admite, además,
\hlpink{una presentación mediante un esquema}.
\end{ptcb}

La exposición sigue a Fantechi~\cite{fantechi}. Nuestros ejemplos principales
son un esquema visto como la categoría \(\Sch/S\), y la categoría
\(\Curves\) de familias de curvas. Primero explicaremos las condiciones de
stack, después añadiremos la algebraicidad y finalmente relacionaremos estas
construcciones con los problemas de módulos. La verificación para un esquema,
incluido el espacio proyectivo, está en la Sección~\ref{sec:scheme-stack}.

\begin{ptcb}
\textbf{Una ruta para seguir la charla sin leer todas las fórmulas.}
Los recuadros de intuición ofrecen una explicación que se puede contar en voz
alta; las definiciones y fórmulas precisan esa explicación. Como primera
imagen, pueden pensar en figuras descritas por ecuaciones y en cómo esas
figuras cambian al variar ciertos parámetros. Un isomorfismo identifica dos
descripciones conservando la estructura que estamos estudiando; no significa
que las descripciones sean literalmente iguales.
\end{ptcb}

\section{¿Qué objetos y qué cubrimientos?}
\label{sec:site}

Una categoría tiene objetos, morfismos, identidades y una composición
asociativa. Por ejemplo, \(\Sch\) tiene esquemas como objetos y morfismos de
esquemas como flechas. Un funtor actúa sobre objetos y flechas, conservando
las identidades y la composición. Durante la charla, los esquemas están
definidos sobre \(\mathbb C\), y todos los morfismos y productos fibrados
respetan esa base.

\begin{ptcb}
\textbf{Intuición: un catálogo con relaciones.}
Pensemos en un catálogo de objetos junto con las comparaciones permitidas
entre ellos. Una categoría registra también cómo encadenar esas comparaciones.
Un funtor traduce el catálogo a otro, respetando ese encadenamiento.
En esta charla, ``esquema'' nombra el tipo de espacio geométrico con el que
trabajamos; como imagen inicial, pueden imaginar figuras descritas por
ecuaciones polinómicas y modelos locales que se pegan.
\end{ptcb}

\begin{ptcb}
Coloquialmente, una \textbf{topología de Grothendieck} especifica qué familias
\(\{S_i\to S\}\) cuentan como cubrimientos. Así, la frase ``estos mapas
cubren \(S\)'' tiene un significado determinado. Los cubrimientos deben ser
compatibles con el cambio de base y con la composición de cubrimientos.
Una categoría con una topología de este tipo se llama \textbf{sitio}.
\end{ptcb}

En la topología de Zariski, los mapas son inmersiones abiertas cuyas imágenes
cubren \(S\). En la \textbf{topología \'etale}, son morfismos \'etale con
imágenes conjuntamente sobreyectivas. No tienen que ser inmersiones abiertas.

Como intuición, entre esquemas complejos lisos, un mapa \'etale tiene
diferencial invertible en cada punto complejo:
\[
  d f_x:T_xU\xrightarrow{\sim}T_{f(x)}S.
\]
Es el análogo algebraico de un difeomorfismo local. Esta descripción mediante
el diferencial se refiere a esquemas lisos; en general, un morfismo \'etale
es un morfismo liso de dimensión relativa cero. Los mapas lisos desempeñan
el papel de las submersiones. Véase~\cite{sp:etale,sp:covers}.

Para un cubrimiento \(u_i:S_i\to S\), los solapamientos son productos
fibrados:
\[
  S_{ij}=S_i\times_S S_j,
  \qquad S_{ijk}=S_i\times_S S_j\times_S S_k.
\]
Para inmersiones abiertas, son las intersecciones habituales. Para
cubrimientos \'etale generales, los productos fibrados sustituyen a esas
intersecciones.

\begin{ptcb}
\textbf{Intuición: varias hojas de un mapa.}
Podemos estudiar un territorio usando varias hojas de un mapa. Un lugar
puede aparecer en más de una hoja, por lo que debemos comparar sus
descripciones. El cubrimiento garantiza que las hojas alcanzan todo el
territorio. En geometría algebraica ampliamos qué puede contar como una
``hoja'': no siempre tiene que ser un simple subconjunto de la base.
\end{ptcb}

\section{Categorías sobre esquemas}
\label{sec:over}

\begin{Def}
Una \textbf{categoría sobre \(\Sch\)} es una categoría \(X\) con un funtor
covariante fijo \(p:X\to\Sch\). Un objeto \(E\) está sobre \(S\) si
\(p(E)=S\), y una flecha \(a\) está sobre \(f\) si \(p(a)=f\).
La \textbf{fibra} \(X_S\) contiene los objetos sobre \(S\) y las flechas
sobre \(\id_S\).
\end{Def}

Una sola categoría registra familias sobre muchos esquemas de parámetros.
La proyección recuerda la base de cada familia y el mapa entre las bases
de cada flecha. La fibra sobre un \(S\) fijo compara familias sobre esa
base particular.

\begin{Ex}[Un esquema como categoría sobre esquemas]
Fijemos un esquema \(S\). Los objetos de \(\Sch/S\) son morfismos
\(f:T\to S\). Una flecha de \(f:T\to S\) a \(f':T'\to S\) es un
morfismo de esquemas \(h:T\to T'\) tal que \(f=f'\circ h\):
\[
\begin{tikzcd}
 T \arrow[rr,"h"] \arrow[dr,"f"'] && T' \arrow[dl,"f'"] \\
 & S
\end{tikzcd}
\]
Su proyección es el \textbf{funtor origen}
\[
 \operatorname{source}_S:\Sch/S\longrightarrow\Sch,
 \qquad (f:T\to S)\longmapsto T,\qquad h\longmapsto h.
\]
En particular, \(f:T\to S\), visto como objeto de esta categoría, está
sobre \(T\). El destino fijo \(S\) es parte de la construcción de la
categoría de objetos sobre \(S\).
\end{Ex}

\begin{Ex}[Familias de curvas]
Un objeto de \(\Curves\) es un morfismo \(\pi:E\to S\), donde \(S\)
es un esquema, \(E\) es un espacio algebraico y \(\pi\) es plano,
propio, de presentación finita y de dimensión relativa a lo sumo uno.
Una flecha es un diagrama cartesiano
\[
\begin{tikzcd}
 E_T \arrow[r,"\bar f"] \arrow[d,"\pi_T"'] & E \arrow[d,"\pi"] \\
 T \arrow[r,"f"'] & S .
\end{tikzcd}
\]
Así, \(E_T\cong E\times_S T\), de forma compatible con los mapas del
diagrama. El funtor \(\operatorname{Target}:\Curves\to\Sch\), que
recuerda el destino de la familia, envía la familia a \(S\) y el diagrama
a \(f\). Aquí \(\pi\) es la proyección de una familia, mientras que
\(\operatorname{Target}\) es un funtor entre categorías.
Esta es la categoría amplia de curvas de~\cite{sp:curves}; todavía no
imponemos género fijo, conexidad ni estabilidad.
\end{Ex}

\begin{ptcb}
\textbf{Intuición: un tablero de parámetros.}
La base \(S\) funciona como un tablero de parámetros. Una familia reúne
los objetos que aparecen al variar esos parámetros. Un mapa \(T\to S\)
permite usar un nuevo tablero: el cambio de base selecciona la familia
correspondiente sobre \(T\). Elegir un solo punto del tablero recupera
un objeto individual de la familia.
\end{ptcb}

\section{Categorías fibradas en grupoides}
\label{sec:cfg}

\begin{ptcb}
Ser una categoría sobre \(\Sch\) nos dice sobre qué bases están los
objetos. Ser \textbf{fibrada} garantiza además que podemos efectuar cambios
de base a lo largo de cualquier mapa de bases. Tener \textbf{fibras que son
grupoides} significa que, sobre una base fija, las flechas son isomorfismos.
\end{ptcb}

\begin{Def}
Una categoría \(p:X\to\Sch\) es \textbf{fibrada} si, para todo
\(f:T\to S\) y todo \(E\in X_S\), existe una flecha cartesiana
\[
 f_E:f^*E\longrightarrow E
 \qquad\text{sobre }f.
\]
Aquí \textbf{cartesiana} expresa la siguiente propiedad universal.
Si \(D\in X_R\), \(g:R\to T\) y \(a:D\to E\) está sobre
\(f\circ g\), existe un único \(b:D\to f^*E\) sobre \(g\) con
\(f_E\circ b=a\). Una \textbf{categoría fibrada en grupoides} (CFG) es
una categoría fibrada tal que cada fibra \(X_S\) es un grupoide.
\end{Def}

Este es el significado categórico general de un levantamiento cartesiano.
En nuestros ejemplos de familias, el levantamiento se representa mediante
el cuadrado usual de cambio de base. En una CFG, toda flecha es cartesiana;
esto se deduce de la condición de grupoide y no forma parte de la
definición de una categoría arbitraria sobre \(\Sch\).
Los cambios de base elegidos son únicos salvo un único isomorfismo
compatible con sus mapas a \(E\). En general,
\(g^*f^*E\cong(f\circ g)^*E\), en lugar de una igualdad literal.

\begin{ptcb}
\textbf{Intuición: cambiar parámetros e identificar descripciones.}
Queremos poder cambiar siempre de tablero de parámetros. El resultado queda
determinado salvo una identificación única que respeta los mapas de la
construcción. La condición de grupoide dice que, sobre una base fija, las
comparaciones admitidas son reversibles. Son identificaciones que conservan
la estructura del objeto.
\end{ptcb}

\begin{Ex}[Los dos ejemplos principales son CFG]
Para \(\Sch/S\), el cambio de base por \(u:T'\to T\) lleva el objeto
\(f:T\to S\) a \(f\circ u:T'\to S\). La flecha correspondiente está
representada por \(u\). La fibra sobre \(T\) es discreta: una flecha sobre
\(\id_T\) existe únicamente entre mapas iguales \(T\to S\), y entonces
es la identidad.

Para \(\Curves\), el cambio de base de \(E\to S\) es
\(E\times_S T\to T\). Las condiciones que definen una familia se
conservan bajo cambio de base. Sobre \(\id_S\), el cuadrado cartesiano
identifica su flecha superior con un isomorfismo de familias. Por tanto,
sus fibras son grupoides.
\end{Ex}

\section{Datos de descenso y efectividad}
\label{sec:descent}

Fijemos una CFG \(p:X\to\Sch\) y un cubrimiento
\(\{u_i:S_i\to S\}\). Escribimos \(E|_T\) para un cambio de base
elegido de \(E\) a \(T\). Todas las comparaciones entre cambios de base
sucesivos usan sus identificaciones canónicas.

\begin{Def}
Un \textbf{dato de descenso} consiste en objetos locales
\(E_i\in X_{S_i}\) e isomorfismos
\[
 \alpha_{ij}:E_i|_{S_{ij}}\xrightarrow{\sim}E_j|_{S_{ij}}
\]
sobre \(\id_{S_{ij}}\), que satisfacen la \textbf{condición de cociclo}
\[
 \alpha_{jk}|_{S_{ijk}}\circ\alpha_{ij}|_{S_{ijk}}
   =\alpha_{ik}|_{S_{ijk}}.
\]
\end{Def}

\begin{ptcb}
Los objetos locales no tienen que ser iguales en los solapamientos.
Disponemos de isomorfismos especificados entre sus cambios de base.
La condición de cociclo exige que dos identificaciones sucesivas coincidan
con la identificación directa.
\end{ptcb}

\begin{ptcb}
\textbf{Intuición: piezas e instrucciones de ensamblaje.}
Tenemos piezas locales e instrucciones precisas para identificar sus zonas
comunes. Con tres piezas, pasar de la primera a la segunda y luego a la
tercera debe coincidir con ir directamente de la primera a la tercera.
La efectividad pregunta si las piezas y esas instrucciones provienen de
una construcción global.
\end{ptcb}

\begin{figure}[htbp]
\centering
\TalkDDDiagram
\caption{Datos de descenso: objetos locales e identificaciones sobre los
solapamientos. Todavía no se ha proporcionado un objeto global sobre \(S\).}
\label{fig:descent-data}
\end{figure}

En un solapamiento triple, la condición de cociclo es la conmutatividad
del triángulo
\[
\begin{tikzcd}[column sep=huge]
 E_i|_{S_{ijk}} \arrow[rr,"\alpha_{ij}"]
   \arrow[dr,"\alpha_{ik}"'] && E_j|_{S_{ijk}} \arrow[dl,"\alpha_{jk}"] \\
 & E_k|_{S_{ijk}} .
\end{tikzcd}
\]
En las etiquetas de las flechas de este triángulo omitimos las restricciones.

\begin{Def}
El dato es \textbf{efectivo} si existen un objeto \(E\in X_S\) e
isomorfismos \(\alpha_i:E|_{S_i}\xrightarrow{\sim}E_i\) tales que
\begin{equation}\label{eq:effective-compatibility}
 \alpha_{ij}
   =(\alpha_j|_{S_{ij}})\circ(\alpha_i|_{S_{ij}})^{-1}.
\end{equation}
\end{Def}

La compatibilidad es esencial: el objeto global debe reproducir las
identificaciones \emph{dadas}; no basta con que sus cambios de base sean
localmente isomorfos a los objetos originales. El recorrido de
\eqref{eq:effective-compatibility} va de \(E_i|_{S_{ij}}\) hacia
\(E|_{S_{ij}}\) y después hacia \(E_j|_{S_{ij}}\).

\begin{figure}[htbp]
\centering
\TalkDDEDiagram
\caption{Un dato efectivo: el objeto global y sus identificaciones con los
objetos locales dados. El triángulo azul expresa la compatibilidad.}
\label{fig:effective-descent}
\end{figure}

En las Figuras~\ref{fig:descent-data}--\ref{fig:effective-descent}, las
líneas punteadas indican el funtor de proyección sobre objetos; no son
flechas adicionales de \(X\). Los mapas \(r_i,r_j,q_i,q_j,s_i,s_j\)
son flechas cartesianas de cambio de base, y las flechas azules son
isomorfismos en las fibras. Partir de un objeto global \(E\) siempre
produce un dato efectivo, usando las comparaciones canónicas entre sus
cambios de base. El axioma de stack pide efectividad de \emph{todo} dato de
descenso, incluidos los que se especifican sin un objeto global.

\section{Los isomorfismos forman un haz}
\label{sec:isom-sheaf}

Esta vez ya tenemos dos objetos globales \(E,E'\in X_S\). Supongamos
que también tenemos isomorfismos locales
\[
 \alpha_i:E|_{S_i}\xrightarrow{\sim}E'|_{S_i}
\]
cuyas restricciones son la misma flecha en cada solapamiento:
\[
 \alpha_i|_{S_{ij}}=\alpha_j|_{S_{ij}}:
 E|_{S_{ij}}\longrightarrow E'|_{S_{ij}}.
\]

\begin{Def}
Decimos que \textbf{los isomorfismos forman un haz} si toda colección
compatible de este tipo proviene de un único isomorfismo
\(\alpha:E\xrightarrow{\sim}E'\) sobre \(\id_S\), con
\(\alpha|_{S_i}=\alpha_i\).
\end{Def}

\begin{figure}[htbp]
\centering
\TalkISDiagram
\caption{Los isomorfismos forman un haz. Las flechas azules locales dadas
coinciden en el solapamiento; la flecha global discontinua debe existir
y ser única.}
\label{fig:isom-sheaf}
\end{figure}

\begin{ptcb}
\textbf{Intuición: comparar dos construcciones ya terminadas.}
Ahora las dos construcciones globales ya existen. En cada zona sabemos
cómo identificar una con la otra. Si estas identificaciones coinciden en
las zonas comunes, deben reunirse en una única identificación global.
Estamos pegando las comparaciones entre las construcciones.
\end{ptcb}

\begin{ptcb}
\textbf{Distinguir las dos preguntas de pegado.}
La efectividad pega \emph{objetos locales con identificaciones especificadas}.
La condición de haz pega \emph{isomorfismos locales compatibles entre dos
objetos globales que ya existen}. Saber únicamente que dos objetos son
localmente isomorfos no basta: los isomorfismos locales elegidos deben
coincidir en los solapamientos.
\end{ptcb}

\section{La definición de stack y las primeras verificaciones}
\label{sec:stack}

\begin{Def}
Un \textbf{stack en grupoides} sobre nuestro sitio es una CFG para la cual
los isomorfismos forman un haz y todo dato de descenso es efectivo.
\end{Def}

En esta charla, ``stack'' significa stack en grupoides. La definición tiene
sentido sobre cualquier sitio, no solo sobre esquemas: se sustituye
\(\Sch\) por una categoría \(\mathcal C\) con una topología de
Grothendieck y se usan sus cubrimientos. Véase~\cite{sp:stack}.

\subsection{Un esquema visto como stack}
\label{sec:scheme-stack}

\begin{Prop}
Para todo esquema \(S\), la categoría \(\Sch/S\) es un stack para la
topología \'etale.
\end{Prop}

\begin{proof}
Es una CFG por la Sección~\ref{sec:cfg}. Sobre un esquema de prueba \(T\),
sus objetos son mapas \(f:T\to S\) y su fibra es discreta.
Los isomorfismos locales entre dos de estos objetos afirman que los dos
mapas coinciden localmente. Por tanto, coinciden globalmente, y el
isomorfismo resultante es único. Esto prueba la condición de haz para los
isomorfismos.

Un dato de descenso consiste en mapas \(f_i:T_i\to S\) cuyos cambios de
base a \(T_i\times_T T_j\) coinciden. Estos se pegan de forma única
en un morfismo \(f:T\to S\), lo que hace efectivo el dato.
Para cubrimientos de Zariski, es el pegado usual de morfismos de esquemas.
Para cubrimientos \'etale, usamos el teorema estándar que afirma que los
morfismos hacia un esquema satisfacen descenso~\cite{sp:morphism-descent}.
\end{proof}

\textbf{El espacio proyectivo es un ejemplo inmediato:}
\(\Sch/\mathbb P^n\) es un stack. Sus objetos sobre \(T\) son mapas
\(T\to\mathbb P^n\), y los mapas locales compatibles se pegan.
Así, el ejemplo sugerido por el organizador se verifica usando exactamente
los dos axiomas anteriores.

\subsection{La categoría de curvas es un stack}

La afirmación para \(\Curves\) es un verdadero teorema de descenso.
Las familias locales compatibles de curvas se pegan respetando las
identificaciones dadas en los solapamientos; los isomorfismos compatibles
entre familias se pegan de forma única. Usar espacios algebraicos como
espacios totales permite el pegado necesario. La familia resultante sigue
siendo plana, propia, de presentación finita y de dimensión relativa a lo
sumo uno. Citamos este teorema, sin desarrollar durante la charla la prueba
del descenso efectivo para espacios algebraicos~\cite{sp:curves}.

En cualquier stack, entre dos realizaciones del mismo dato de descenso
existe un único isomorfismo \emph{compatible con las identificaciones
locales elegidas}. La efectividad da la existencia; la condición de haz
para los isomorfismos proporciona esta unicidad al pegar las comparaciones
locales.

\section{Representabilidad y morfismos de stacks}
\label{sec:representability}

Abreviamos el stack \(\Sch/S\) por \(S\) cuando el contexto es claro.
Escribir un esquema como un stack incluye siempre su proyección al origen
hacia \(\Sch\).

\begin{Def}
Siguiendo la convención de Fantechi, un stack \(X\) es
\textbf{representable por un esquema} si
\(X\simeq\Sch/S\) para algún esquema \(S\). Aquí \(\simeq\) significa
una equivalencia sobre \(\Sch\), compatible con la estructura categórica.
\end{Def}

Un \textbf{morfismo de stacks} \(F:X\to Y\) es un funtor sobre la base:
\(p_Y\circ F=p_X\). Esta condición se aplica tanto a las flechas como
a los objetos. Es un funtor entre las categorías de familias.

Por ejemplo, un morfismo \(\Sch/S\to X\) determina el objeto
\(F(\id_S)\in X_S\). Recíprocamente, un objeto \(E\in X_S\) da el
funtor que envía \(f:T\to S\) a \(f^*E\), con la acción inducida sobre
las flechas. Estas construcciones son inversas salvo los isomorfismos
naturales correspondientes.

\begin{Def}
Un morfismo \(F:X\to Y\) es \textbf{representable por esquemas} si, para
todo esquema \(T\) y todo morfismo \(G:\Sch/T\to Y\), el producto
fibrado \((\Sch/T)\times_Y X\) está representado por un esquema \(Z_G\).
\end{Def}

Así, examinamos \(F\) mediante el siguiente diagrama:
\begin{equation}\label{eq:test-representability}
\begin{tikzcd}[column sep=3.5em]
 Z_G \arrow[d,"\rho_G"',""{name=scheme-column,anchor=center}]
 & (\Sch/T)\times_Y X
     \arrow[d,"\pr"',""{name=stack-column,anchor=center}] \arrow[r]
 & X \arrow[d,"F"] \\
 T & \Sch/T \arrow[r,"G"'] & Y
 \arrow[from=scheme-column,to=stack-column,dashed,
   shorten <=.5em,shorten >=1em,
   "\text{representa}"{xshift=-.25em}]
\end{tikzcd}
\end{equation}
Solo el cuadrado derecho es el producto fibrado de stacks; la flecha
discontinua explica cómo el morfismo de esquemas de la izquierda representa
su proyección. El esquema \(Z_G\) y el morfismo \(\rho_G\) dependen de
\(F\) y del morfismo de prueba \(G\). No se exige que \(X\) o \(Y\)
sean representables.

\begin{ptcb}
\textbf{Intuición: examinar con parámetros conocidos.}
Aunque \(X\) e \(Y\) sean stacks, podemos mirar \(F\) usando una
familia parametrizada por un esquema ordinario \(T\). La representabilidad
dice que, después de esta prueba, obtenemos de nuevo un esquema
\(Z_G\). Esto permite estudiar el morfismo usando geometría de esquemas.
\end{ptcb}

El producto fibrado cumple el papel esperado de un cambio de base, con la
conmutatividad entendida salvo un isomorfismo natural especificado.
Concretamente, sobre \(T\), un objeto de \(X\times_Y X'\) consiste en
objetos \(E\in X_T\), \(E'\in X'_T\), junto con un isomorfismo entre
sus imágenes en \(Y_T\). Este isomorfismo adicional permite que el
producto fibrado conserve los automorfismos. Esta descripción basta aquí,
sin desarrollar todo el formalismo de las \(2\)-categorías.

Si \(P\) es una propiedad de los morfismos de esquemas que se conserva
por cambio de base, un morfismo representable \(F\) tiene \(P\) cuando
todo \(\rho_G\) en \eqref{eq:test-representability} tiene \(P\). Así se
definen la propiedad de ser liso, la propiedad de ser étale y la sobreyectividad de
nuestros morfismos representables.

\begin{ptcb}
\textbf{Sobreyectividad de un morfismo de stacks.}
La condición es que todo morfismo de esquemas de prueba
\(\rho_G:Z_G\to T\) sea sobreyectivo. El funtor subyacente no tiene por
qué alcanzar todos los objetos globales.
Para un atlas liso o étale y sobreyectivo, todo objeto proviene del atlas
\emph{localmente}, salvo isomorfismo. Por ejemplo, el atlas
\(\operatorname{Spec}\mathbb C\to V_r\) proporciona fibrados triviales,
y todo fibrado vectorial de rango \(r\) es localmente trivial.
\end{ptcb}

\section{Presentaciones: Artin y Deligne--Mumford}
\label{sec:algebraic}

\begin{Def}
Un stack \(X\) es \textbf{algebraico en el sentido de Artin} si existe
un esquema \(U\) y un morfismo representable, \emph{liso y sobreyectivo}
\(U\to X\). Es \textbf{Deligne--Mumford} (DM) si puede elegirse tal
morfismo \emph{étale y sobreyectivo}. Un morfismo así es una
\textbf{presentación}, también llamada \textbf{atlas} liso o étale.
\end{Def}

Todo stack DM es Artin, pues los morfismos étale son lisos. Un morfismo
representable \(U\to X\), por sí solo, no es una presentación: son
esenciales la condición de propiedad de ser liso o de ser étale y la
\emph{sobreyectividad}.

La terminología de atlas recuerda al atlas de una variedad diferencial:
las inversas de las cartas coordenadas
\(V_a\subseteq\mathbb R^n\to M\) se reúnen en un difeomorfismo local
sobreyectivo \(\coprod_a V_a\to M\). Un atlas étale cumple un papel
similar. Sus intersecciones \(U\times_X U\) recuerdan isomorfismos entre
las descripciones en cartas, incluidos los automorfismos, y no solamente
intersecciones de subconjuntos abiertos.

\begin{nonum-Rmk}[Convención para el teorema de las curvas]
La definición moderna permite que ``representable'' signifique
representable por \emph{espacios algebraicos}; también incluye
explícitamente la representabilidad de la diagonal. Un atlas representable,
liso y sobreyectivo proporciona esa condición sobre la diagonal. Fantechi
usa la convención más restrictiva de esquemas y restringe la categoría
base. Usamos la convención habitual de espacios algebraicos al citar el
teorema de que el stack completo \(\Curves\) es Artin, y permitimos
esquemas complejos generales como bases. Las formulaciones étale y fppf
dan los mismos stacks algebraicos.
Referencias:~\cite{sp:algebraic,sp:atlas,sp:etale-fppf}.
\end{nonum-Rmk}

\begin{Th}[El criterio de la diagonal]
Para un stack \(X\) que ya sabemos que es Artin,
\[
  X\text{ es DM}
  \quad\Longleftrightarrow\quad
  \Delta_X:X\longrightarrow X\times X\text{ es no ramificada}.
\]
\end{Th}

Aquí la diagonal envía un objeto al par formado por ese mismo objeto dos
veces. Su cambio de base a lo largo de un par \(E,E'\in X_T\)
parametriza isomorfismos de \(E\) a \(E'\). En particular, tomando
\(E=E'\), obtenemos su grupo de automorfismos. El criterio es un teorema,
no un axioma adicional que deba comprobarse después de encontrar un atlas
étale~\cite{sp:dm}.

Equivalentemente, los estabilizadores geométricos deben ser no
ramificados. Sobre \(\mathbb C\), si estos estabilizadores son de tipo
finito, esto equivale a que todo objeto geométrico tenga un número finito
de automorfismos. Estas hipótesis se cumplen en los ejemplos de curvas
que veremos. Así, ``automorfismos finitos'' es un criterio \emph{entre
stacks que ya sabemos que son algebraicos}, y no una definición que por
sí sola produzca algebraicidad~\cite{sp:dm-locus,sp:curve-aut}.

\begin{ptcb}
\textbf{Intuición: simetrías finitas frente a simetrías continuas.}
Como imagen intuitiva, un triángulo equilátero tiene un número finito de
simetrías rígidas, mientras que un círculo admite una familia continua
de rotaciones. Esta es solo una comparación geométrica: el criterio DM
se refiere a los estabilizadores algebraicos y requiere las hipótesis
anteriores. En los ejemplos sobre \(\mathbb C\), ayuda a distinguir
simetrías finitas de familias continuas de automorfismos.
\end{ptcb}

\subsection{Retomamos los ejemplos}

\begin{Ex}[Los esquemas ya son stacks DM]
Para todo esquema \(S\), la presentación identidad
\(\Sch/S\to\Sch/S\) es representable, étale y sobreyectiva.
Por tanto, \(\Sch/S\) es DM y Artin, incluso cuando \(S\) es singular.
Esta afirmación se refiere al atlas identidad, y no a la propiedad de ser liso de
\(S\to\operatorname{Spec}\mathbb C\).
Para el espacio proyectivo, sus cartas afines habituales dan el atlas
étale más visual \(\coprod_{i=0}^n\mathbb A^n\to\mathbb P^n\).
\end{Ex}

\begin{Th}[El stack de curvas]
El stack \(\Curves\) es Artin. En general, no es DM.
\end{Th}

La algebraicidad es un teorema sustancial~\cite{sp:curves}. Para ver que
la condición DM puede fallar, consideremos el objeto
\(\mathbb P^1\to\operatorname{Spec}\mathbb C\) de \(\Curves\).
Su grupo de automorfismos es \(\operatorname{PGL}_2\), de dimensión
tres, por lo que no cumple la condición de no ramificación sobre
\(\mathbb C\)~\cite{sp:curve-aut}.

\begin{ptcb}
El stack representable \(\Sch/\mathbb P^1\) es DM.
La curva \(\mathbb P^1\), como \emph{objeto de \(\Curves\)}, tiene
automorfismos continuos no triviales. Estas afirmaciones se refieren a
categorías diferentes: la primera clasifica morfismos \emph{hacia}
\(\mathbb P^1\); la segunda categoría clasifica familias de curvas.
\end{ptcb}

Para el ejemplo estable con puntos marcados, fijemos \(g,n\geq0\) con
\(2g-2+n>0\). La categoría \(\overline{\mathcal M}_{g,n}\) tiene como
objetos familias de curvas como las anteriores, equipadas con marcas
\((E\to S;\sigma_1,\ldots,\sigma_n)\), cuyas fibras geométricas son
conexas y nodales, de género aritmético \(g\), con marcas etiquetadas y
disjuntas en el lugar liso, y que satisfacen la condición de estabilidad
\[
  2g(D)-2+m_D>0
\]
en cada componente \(D\) de la normalización de una fibra geométrica;
\(m_D\) cuenta las marcas y las ramas de nodos sobre esa componente.
También se permiten fibras lisas, sin nodos.
Los morfismos son diagramas cartesianos que conservan las marcas.

\begin{Th}[Curvas estables con puntos marcados]
La categoría \(\overline{\mathcal M}_{g,n}\) es un stack DM.
\end{Th}

El teorema incluye la algebraicidad. La estabilidad hace que el grupo de
automorfismos de toda curva geométrica con puntos marcados sea finito;
los automorfismos deben conservar \emph{cada marca etiquetada}.
Véase~\cite[\S6.3 y Teorema~6.4.16]{alper}.
Para \(n>0\), este es un problema de módulos con puntos marcados:
olvidar las marcas da un morfismo hacia \(\Curves\), pero la categoría
con marcas no es literalmente una subcategoría con los mismos objetos
que \(\Curves\).

\begin{center}
\begin{tabular}{lll}
\toprule
Stack & ¿Artin? & ¿DM? \\
\midrule
\(\Sch/S\) & Sí & Sí \\
\(\Curves\) & Sí & En general, no \\
\(\overline{\mathcal M}_{g,n}\), \(2g-2+n>0\) & Sí & Sí \\
\bottomrule
\end{tabular}
\end{center}

\section{Si queda tiempo: problemas de módulos y espacio proyectivo}
\label{sec:moduli}

Un \textbf{problema de módulos} busca clasificar un tipo especificado de
objetos geométricos, con sus condiciones y los isomorfismos permitidos,
y describir sus familias de forma compatible con el cambio de base.
La categoría de todas esas familias sobre \(\Sch\) codifica el problema.
Probar el descenso y la algebraicidad le da la estructura de un stack
algebraico de módulos.

\begin{ptcb}
\textbf{Intuición: clasificar formas y permitir que varíen.}
Primero decidimos qué objetos estudiamos y cuándo dos descripciones
representan el mismo objeto: por ejemplo, mediante un isomorfismo que
respete las marcas. Después preguntamos cómo pueden variar esos objetos
en familias. El problema de módulos reúne la clasificación y esa variación,
junto con las simetrías de cada objeto.
\end{ptcb}

Para curvas estables con puntos marcados, la pregunta es: fijados \(g,n\),
clasificar las curvas estables de género \(g\) con \(n\) puntos lisos
marcados y etiquetados, salvo isomorfismos que conserven las marcas,
y entender cómo varían en familias.
El stack \(X=\overline{\mathcal M}_{g,n}\) conserva las familias, sus
isomorfismos y sus relaciones de cambio de base. Un morfismo de stacks
\(S\to X\) es una familia sobre \(S\); en la descripción mediante
funtores, esa familia es la imagen de \(\id_S\).

Al pasar a clases de isomorfismo obtenemos un \textbf{funtor de módulos}
con valores en conjuntos:
\[
 F_X:\Sch^{\mathrm{op}}\longrightarrow\mathbf{Sets},
 \qquad S\longmapsto\{\text{objetos de }X_S\}/\cong,
\]
donde \(u:T\to S\) actúa mediante \([E]\mapsto[u^*E]\).
Este funtor contravariante es diferente de la proyección covariante
\(p:X\to\Sch\). Conservar los grupoides \(X_S\), en lugar de solo sus
clases de isomorfismo, permite conservar también los automorfismos.

Un \textbf{espacio de módulos fino} para el funtor con valores en conjuntos
es un esquema \(M\) con biyecciones naturales
\[
 F_X(S)\cong\operatorname{Hom}_{\Sch}(S,M)
\]
para todas las bases \(S\), compatibles con el cambio de base. No basta con
hacer corresponder las clases de isomorfismo sobre
\(\operatorname{Spec}\mathbb C\). Un espacio de módulos grueso satisface
una propiedad universal más débil y clasifica las clases de isomorfismo
geométricas; no tiene por qué clasificar familias arbitrarias de esta
manera. Escribir \(\overline{\mathcal M}_{g,n}\) para el stack y
\(\overline M_{g,n}\) para su espacio de módulos grueso mantiene visible
la distinción.

\subsection{El problema de módulos concreto representado por el espacio proyectivo}

Además de la comprobación directa de la condición de stack en la
Sección~\ref{sec:scheme-stack}, el espacio proyectivo tiene la siguiente
interpretación de módulos:
\[
 \operatorname{Hom}(T,\mathbb P^n)
 \cong
 \left\{
   \mathcal O_T^{\oplus(n+1)}\twoheadrightarrow L
   \;:\;L\text{ es un fibrado en rectas}
 \right\}\big/\cong.
\]
Aquí los isomorfismos conmutan con los morfismos cociente. Equivalentemente,
se especifica un fibrado en rectas con \(n+1\) secciones globales
ordenadas que lo generan. Estos datos se transportan por cambio de base a lo
largo de morfismos de esquemas de parámetros~\cite{sp:projective}.

Un automorfismo de \(L\) que conserve el cociente debe ser la identidad,
pues las secciones generadoras lo determinan. Junto con la biyección
natural mostrada, esto prueba que el problema con valores en grupoides
está representado por \(\Sch/\mathbb P^n\). Si, en cambio, clasificamos
\emph{todos} los fibrados en rectas, sin generadores elegidos ni condición
de generación, obtenemos el stack conocido \(V_1\simeq B\mathbb G_m\).
La multiplicación por un escalar no nulo ya es un automorfismo sobre
\(\operatorname{Spec}\mathbb C\).

\begin{ptcb}
\textbf{Posible frase de cierre.}
Un esquema puede verse como un stack de morfismos hacia él. Un stack de
módulos permite hacer geometría con familias y sus isomorfismos cuando
un esquema ordinario de parámetros no conserva la información que
necesitamos.
\end{ptcb}

\begin{thebibliography}{SP10}
\bibitem[F]{fantechi}
Barbara Fantechi, \emph{Stacks for everybody}, European Congress of
Mathematics, 2001. Fuente principal para la progresión de definiciones;
el presente texto es un borrador independiente para una exposición.

\bibitem[SP1]{sp:etale}
The Stacks Project Authors, \emph{\'Etale morphisms} y
\emph{\'Etale and smooth morphisms}.
\href{https://stacks.math.columbia.edu/tag/02GH}{Tag 02GH};
\href{https://stacks.math.columbia.edu/tag/039P}{Tag 039P}.

\bibitem[SP2]{sp:covers}
The Stacks Project Authors, \emph{The \'etale topology}.
\href{https://stacks.math.columbia.edu/tag/0214}{Tag 0214}.

\bibitem[SP3]{sp:stack}
The Stacks Project Authors, \emph{Stacks in groupoids}, Definición~8.5.1.
\href{https://stacks.math.columbia.edu/tag/02ZH}{Tag 02ZH}.

\bibitem[SP4]{sp:morphism-descent}
The Stacks Project Authors, Lema~35.13.7: los funtores representables
satisfacen el descenso para cubrimientos fpqc y, por tanto, para
cubrimientos étale.
\href{https://stacks.math.columbia.edu/tag/023Q}{Tag 023Q}.

\bibitem[SP5]{sp:curves}
The Stacks Project Authors, \emph{The stack of curves},
Situación~99.15.1, los lemas de CFG y descenso, y Teorema~99.15.11.
\href{https://stacks.math.columbia.edu/tag/0D4Y}{Tag 0D4Y}.

\bibitem[SP6]{sp:algebraic}
The Stacks Project Authors, definiciones de stacks algebraicos y
Deligne--Mumford.
\href{https://stacks.math.columbia.edu/tag/026O}{Tag 026O};
\href{https://stacks.math.columbia.edu/tag/03YO}{Tag 03YO}.

\bibitem[SP7]{sp:atlas}
The Stacks Project Authors, Lema~94.15.3: algebraicidad a partir de un
cubrimiento representable, liso y sobreyectivo.
\href{https://stacks.math.columbia.edu/tag/05UL}{Tag 05UL}.

\bibitem[SP8]{sp:etale-fppf}
The Stacks Project Authors, \emph{Algebraic stacks in the \'etale topology},
Lema~97.19.1.
\href{https://stacks.math.columbia.edu/tag/076U}{Tag 076U}.

\bibitem[SP9]{sp:dm}
The Stacks Project Authors, Definición~101.4.1 y Teorema~101.21.6:
diagonal no ramificada y presentaciones étale.
\href{https://stacks.math.columbia.edu/tag/04YW}{Tag 04YW};
\href{https://stacks.math.columbia.edu/tag/06N3}{Tag 06N3}.

\bibitem[SP10]{sp:dm-locus}
The Stacks Project Authors, Lema~101.22.2: el lugar DM y los
estabilizadores geométricos.
\href{https://stacks.math.columbia.edu/tag/0DSL}{Tag 0DSL}.

\bibitem[SP11]{sp:curve-aut}
The Stacks Project Authors, \emph{Curves with finite reduced automorphism
groups}, en particular el Ejemplo~109.7.2 y el Lema~109.7.3.
\href{https://stacks.math.columbia.edu/tag/0DST}{Tag 0DST}.

\bibitem[SP12]{sp:projective}
The Stacks Project Authors, Lemas~27.13.1 y 27.13.3: la interpretación de
módulos y las cartas afines habituales del espacio proyectivo.
\href{https://stacks.math.columbia.edu/tag/01NE}{Tag 01NE};
\href{https://stacks.math.columbia.edu/tag/01NG}{Tag 01NG}.

\bibitem[A]{alper}
Jarod Alper, \emph{Stacks and Moduli}, Capítulo~6, en particular \S6.3 y
Teorema~6.4.16 (curvas estables con puntos marcados).
\href{https://sites.math.washington.edu/~jarod/moduli.pdf}{Notas del autor}.
\end{thebibliography}

\end{document}
